\chapter{Differential forms on Berkovich spaces}
\index{differential form!CLD}\index{Berkovich space!differential forms}
\label{chap:cld-forms}

The preceding chapters built the real differential calculus in three stages:
superforms on affine spaces, forms on calibrated tropical spaces, and
tropicalizations of analytic spaces.  We now assemble these stages into the
Chambert--Loir--Ducros sheaf of forms on a Berkovich space.  The main point is
not a new coordinate formula.  It is the proof that the chartwise formulae are
independent of the moment map and therefore belong intrinsically to the
analytic space.

It is worth stating at the outset exactly what has to be done, because the
construction looks deceptively short.  \Cref{chap:analytic-tropicalization}
produced, for every moment map $f$ on an analytic domain $V$ of $X$, a
continuous map $f_{\mathrm{trop}}:V\to N_{\mathbf R}$ with polyhedral image,
a distinguished locus $\Sigma_f\subset V$ where the coordinates attain
maximal rank, and a canonical calibration on $\Sigma_f^{(n)}$ that does not
depend on $f$.  \Cref{chap:superforms,chap:tropical-spaces,chap:tropical-integration}
produced a complete differential calculus --- wedge products, $\dprime$,
$\ddprime$, positivity cones, integration, Stokes' formula --- on polyhedral
supports.  The obvious move is to declare a form on $X$ to be, locally, the
pullback of a superform along a tropicalization.  Three things then have to
be proved, and none of them is formal.

First, \emph{the pullback notation is a fiction}.  The map
$f_{\mathrm{trop}}$ is only continuous; $X$ is not a manifold; there is no
differential of $f_{\mathrm{trop}}$ and no classical pullback of forms.  What
the notation abbreviates is a class in a filtered colimit taken over the
category of tropical charts.  Everything therefore depends on knowing exactly
when two chart representatives define the same object, that is, on computing
the kernel of $\alpha\mapsto f_{\mathrm{trop}}^*\alpha$.  That computation is
\Cref{prop:ch13-nullity}, and its polyhedral engine is
\Cref{lem:ch13-surjective-pullback-injective}: a superform on a polyhedron
is determined by its pullback along any surjective integer-affine map onto
it.

Second, \emph{the resulting presheaf is not a sheaf, and is not even
separated}.  The failure is not pathological bookkeeping: it is visible on a
discrete tropical space, as the warning following \Cref{prop:ch13-nullity}
records.  One must sheafify, and one must then check that the sheafification
is still described by charts.  This is \Cref{thm:ch13-sheaf}.

Third, \emph{the invariant that integration needs is not the tropical image
but the calibration}.  Two charts see the same analytic space through
different polyhedra of possibly different dimensions, related only through
their product.  The reason $\int_X\omega$ is well defined is not that the
polyhedra agree --- they do not --- but that the canonical calibrations of
\Cref{def:trop-canonical-calibration} agree, which is
\Cref{thm:trop-calibration-chart-independent}.  All of
\Cref{sec:ch13-integration,sec:ch13-stokes} is an unfolding of that one
theorem.

Once these three points are settled, the chapter's structural statement is
that a topologically separated $K$-analytic space purely of dimension $n$,
equipped with its analytic boundary, \emph{is} a calibrated tropical space of
level $n$ (\Cref{thm:analytic-is-tropical-space}).  Stokes' formula on $X$ is
then not a new theorem but the tropical Stokes formula
\Cref{thm:cld-global-stokes} read through that identification, and the
substance lies in identifying the analytic boundary with the tropical one.

\paragraph{Roadmap.}
\Cref{sec:ch13-charts} constructs the sheaves $A_X^{p,q}$, proves the
nullity criterion and the sheaf property, compares ordinary forms with
$G$-forms, and computes the first examples on the Berkovich affine line and
on an annulus.  \Cref{sec:ch13-operations} descends the algebraic and
differential operations, proves their locality, constructs the pullback along
an analytic morphism, and then devotes a subsection to the question of what
that pullback is and is not.  \Cref{sec:ch13-topology} fixes the topology on
test forms; \Cref{sec:ch13-positivity} the positivity cones.
\Cref{sec:ch13-integration} integrates top-degree forms on a pure-dimensional
space, deducing well-definedness from the chart independence of
\Cref{chap:analytic-tropicalization}, and identifies $X$ as a tropical space.
\Cref{sec:ch13-stokes} states and proves Stokes' formula on analytic domains
with boundary, with the hypotheses spelled out, and verifies it on an annulus
with explicit numbers.  \Cref{sec:ch13-skeleton} compares forms on $X$ with
forms on a skeleton and computes the Tate curve completely.
\Cref{sec:ch13-properness} isolates where properness and boundarylessness are
genuinely required.  \Cref{sec:ch13-partitions} records the partition of
unity theorem in the exact form used later, and
\Cref{sec:ch13-dictionary} collects the dictionary.

\paragraph{Dependencies.}
From \Cref{chap:superforms}: the superform algebra, the operators of
\Cref{def:ch09-operators}, the identities of
\Cref{prop:superform-differential-identities}, contraction
(\Cref{def:ch09-contraction}), affine pullback
(\Cref{prop:ch09-pullback-functorial}), and the positivity cones of
\Cref{def:ch09-cones}.  From \Cref{chap:tropical-spaces}: forms on a
polyhedral support (\Cref{def:ch10-forms-on-support}) with their kernel
\eqref{eq:ch10-form-kernel}, the sheaf property
(\Cref{lem:ch10-forms-sheaf}), presentation independence
(\Cref{thm:ch10-forms-presentation-independent}) and its product form
(\Cref{cor:ch10-forms-graph}), common refinements
(\Cref{prop:tropical-common-refinement}), and calibrations
(\Cref{def:ch10-calibration}).  From \Cref{chap:tropical-integration}:
lattice-normalized integration (\Cref{def:ch11-lattice-measure}), subdivision
independence (\Cref{prop:integral-subdivision-independent}), boundary traces
\eqref{eq:superform-boundary-traces}, the defect formula
\eqref{eq:ch11-defect-stokes}, and the Stokes and Green theorems
\Cref{thm:balanced-stokes,thm:cld-global-stokes,thm:ch11-green-defect}.  From
\Cref{chap:analytic-tropicalization}: moment maps
(\Cref{def:trop-moment-invariant}), tropical charts and products
(\Cref{lem:product-moment-map,lem:trop-refinement-calculus}), local skeleta
(\Cref{def:trop-local-skeleton,thm:trop-sigma-characterization}), the
canonical calibration and its chart independence
(\Cref{def:trop-canonical-calibration,thm:trop-calibration-chart-independent,cor:trop-calibration-glue}),
tropical polyhedra and multiplicities
(\Cref{def:trop-tropical-polyhedron,thm:trop-multiplicity-formula}), thinness
of boundaries (\Cref{prop:trop-boundary-thin,cor:trop-boundary-negligible}),
and partitions of unity (\Cref{thm:cld-partitions-unity}).

\paragraph{Standing hypotheses.}
Throughout this chapter, $K$ is a complete non-Archimedean field and $X$ is a
topologically separated $K$-analytic space.  When integration is used, $X$ is
purely $n$-dimensional.  We distinguish ordinary open covers from
$G$-covers by analytic domains whenever this distinction matters.
$K$ is not assumed algebraically closed, nontrivially valued, or of residue
characteristic zero; $X$ is not assumed good, reduced, or boundaryless unless
this is said.  Each of goodness, paracompactness, holomorphic separation,
boundarylessness and purity of dimension is a genuine hypothesis somewhere
below, and \Cref{sec:ch13-properness} lists which statement needs which.

\paragraph{What is imported.}
Four analytic facts are used without proof, each with an exact reference: the
comparison of ordinary and $G$-germs on an almost good space
(\cite[Proposition~15.1.3]{CLD25}, used in
\Cref{prop:cld-local-normal-form}); the vanishing of a form of bidegree
$(p,q)$ off the locus $d_K\ge\max(p,q)$
(\cite[Proposition~15.2.1]{CLD25}); the holomorphic separation and
Stone--Weierstrass package
(\cite[Corollary~15.2.5, Theorem~15.2.6, Proposition~15.2.8]{CLD25}, used in
\Cref{thm:cld-smooth-partition}); and the identification of the analytic
boundary with the parlinear boundary
(\cite[Theorems~15.6.1 and~15.7.1]{CLD25}, used in
\Cref{thm:analytic-is-tropical-space}).  Everything else in the chapter is
proved here or reduced to a statement proved in
\Cref{chap:superforms,chap:tropical-spaces,chap:tropical-integration,chap:analytic-tropicalization}.

\section{From tropical charts to a sheaf of forms}
\label{sec:ch13-charts}

\subsection{The presheaf of chart representatives}
\label{sec:ch13-presheaf}

Let $V$ be an analytic domain of $X$ and let
\[
             f=(f_1,\ldots,f_r):V\longrightarrow\mathbf G_{m,K}^r
\]
be a moment map.  We use the convention
\[
 f_{\mathrm{trop}}(x)=(-\log|f_1(x)|,\ldots,-\log|f_r(x)|).
\]
A smooth superform $\alpha$ on a polyhedral neighborhood of
$f_{\mathrm{trop}}(V)$ determines a form denoted
$f_{\mathrm{trop}}^*\alpha$.  As in \Cref{chap:analytic-tropicalization},
this notation denotes a chart representative, not the classical pullback by
an everywhere differentiable map.

We make the bookkeeping explicit, because the whole construction is a
colimit and the reader must know over what.

\begin{definition}[The chart category]
\label{def:ch13-chart-category}
\index{tropical chart!category of}
Let $V$ be an analytic domain of $X$.  A \emph{tropical chart datum} on $V$
is a pair $(f,P)$ consisting of a moment map $f:V\to T$ and a closed
polyhedral subset $P\subset T_{\mathrm{trop}}$ containing
$\Trop_f(V)=f_{\mathrm{trop}}(V)$.  A morphism
$(f,P)\to(f',P')$ is an integer-affine map
$F:T_{\mathrm{trop}}\to T'_{\mathrm{trop}}$ with
\[
 f'_{\mathrm{trop}}=F\circ f_{\mathrm{trop}}
 \quad\text{on }V,
 \qquad F(P)\subset P'.
\]
The resulting category is denoted $\mathscr T(V)$.  We say that $(f,P)$
\emph{dominates} $(f',P')$ when such an $F$ exists; this is the refinement
relation of \Cref{lem:trop-refinement-calculus} with the polyhedral targets
made explicit.
\end{definition}

\begin{lemma}
\label{lem:ch13-charts-filtered}
Let $V$ be an analytic domain of $X$.
\begin{enumerate}
\item\label{itm:ch13-cf:final} $\mathscr T(V)$ has a terminal object, namely
  the constant moment map to the trivial torus with $P$ a point.
\item\label{itm:ch13-cf:prod} Any two objects $(f,P)$ and $(f',P')$ of
  $\mathscr T(V)$ are dominated by a common object, namely
  $\bigl((f,f'),P\times P'\bigr)$, through the two coordinate projections.
\item\label{itm:ch13-cf:equal} Suppose $V$ is compact and let
  $F_1,F_2:(f,P)\to(f',P')$ be two morphisms.  Then $F_1$ and $F_2$ agree on
  $\Trop_f(V)$, and there is an object of $\mathscr T(V)$ admitting a
  morphism to $(f,P)$ on which the pullbacks $F_1^*$ and $F_2^*$ agree.
\end{enumerate}
Consequently, for $V$ compact, the opposite category
$\mathscr T(V)^{\mathrm{op}}$ is filtered, and colimits indexed by it may be
computed with single representatives.
\end{lemma}

\begin{proof}
\cref{itm:ch13-cf:final} The trivial torus $T=\Spec K$ has
$T_{\mathrm{trop}}=\{0\}$, and any $(f,P)$ receives the unique affine map
to $\{0\}$; there is nothing to check.

\cref{itm:ch13-cf:prod} By \Cref{lem:product-moment-map} the pair
$(f,f'):V\to T\times T'$ is a moment map, and by
\Cref{lem:trop-moment-basic}\,\cref{itm:trop-moment:prod} its tropicalization
is $(f_{\mathrm{trop}},f'_{\mathrm{trop}})$.  Hence
$\Trop_{(f,f')}(V)\subset P\times P'$, and the two coordinate projections
are integer-linear maps carrying $P\times P'$ into $P$ and $P'$ and
satisfying the required identities.

\cref{itm:ch13-cf:equal} For $x\in V$,
$F_1(f_{\mathrm{trop}}(x))=f'_{\mathrm{trop}}(x)=F_2(f_{\mathrm{trop}}(x))$,
so $F_1=F_2$ on $\Trop_f(V)$.  This does not yet make $F_1^*=F_2^*$ as maps
$A^{p,q}(P')\to A^{p,q}(P)$, since $P$ may be strictly larger than
$\Trop_f(V)$.  But $V$ is compact, so $\Trop_f(V)$ is itself a compact
polyhedral subset of $T_{\mathrm{trop}}$
(\Cref{thm:analytic-tropical-polyhedrality}), and
$\bigl(f,\Trop_f(V)\bigr)$ is an object of $\mathscr T(V)$ mapping to
$(f,P)$ by the identity of $T_{\mathrm{trop}}$.  On it the
two pullbacks agree.  Indeed, let $x$ lie in a cell $\sigma$ of
$\Trop_f(V)$ and let $v$ belong to the cone of directions along which
$\sigma$ is entered at $x$; that cone spans $N_{\sigma,\mathbf R}$, since
$\sigma$ is convex of full dimension in its affine span.  For such $v$ the
points $x+tv$ lie in $\Trop_f(V)$ for small $t>0$, so the difference
quotients $t^{-1}\bigl(F_i(x+tv)-F_i(x)\bigr)$ are independent of
$i\in\{1,2\}$, and by linearity the linear parts $DF_1$ and $DF_2$ agree on
$N_{\sigma,\mathbf R}$.  Hence $F_1^*\alpha$ and $F_2^*\alpha$ have the same
evaluations on cell tangent vectors at every point of $\Trop_f(V)$, that is,
they differ by an element of the kernel \eqref{eq:ch10-form-kernel}.  This is
the argument of Step~3 of the proof of
\Cref{thm:ch10-forms-presentation-independent}.

Filteredness of the opposite category is exactly
\cref{itm:ch13-cf:prod} and \cref{itm:ch13-cf:equal}.
\end{proof}

\begin{warning}{Compactness enters already in the bookkeeping}
Without compactness of $V$ the tropical image $\Trop_f(V)$ need not be a
polyhedral subset, two morphisms of chart data need not induce the same
transition map, and the indexing category is then a preorder without a
canonical transition system.  This is the mildest appearance of the
phenomenon that becomes the failure of separatedness recorded after
\Cref{prop:ch13-nullity}.  It is harmless in practice, because forms are
defined by local representability and every comparison made below is carried
out on compact analytic domains; but it is the reason the definitions that
follow are phrased through germs rather than through a global colimit.
\end{warning}

\begin{definition}[Chart representatives]
\label{def:ch13-presheaf}
For $p,q\ge0$ and an analytic domain $V$ of $X$, set
\[
 A^{p,q}_{\mathrm{pf}}(V)
 =\varinjlim_{(f,P)\in\mathscr T(V)^{\mathrm{op}}}A^{p,q}(P),
\]
the colimit being taken along the pullback maps $F^*$ and $A^{p,q}(P)$
denoting superforms on the polyhedral support $P$ in the sense of
\Cref{def:ch10-forms-on-support}.  The class of $\alpha\in A^{p,q}(P)$ is
written $f_{\mathrm{trop}}^*\alpha$.  Restriction along an inclusion of
analytic domains makes $V\mapsto A^{p,q}_{\mathrm{pf}}(V)$ a presheaf on the
$G$-topology of $X$, and its restriction to open subsets a presheaf on the
topological space $X$.  For $V$ compact the colimit is filtered by
\Cref{lem:ch13-charts-filtered}; in general we use it only through the
compact case and through germs.
\end{definition}

For $V$ compact, by \Cref{lem:ch13-charts-filtered} the colimit is filtered,
so every element
of $A^{p,q}_{\mathrm{pf}}(V)$ is of the form $f_{\mathrm{trop}}^*\alpha$ for
a single chart, and
$f_{\mathrm{trop}}^*\alpha=g_{\mathrm{trop}}^*\beta$ if and only if the two
pullbacks agree in some chart dominating both --- equivalently, by
\cref{itm:ch13-cf:prod}, in the product chart.  This is precisely the
identification used in the definition below.  What is \emph{not} yet visible
is how large the kernel of $\alpha\mapsto f_{\mathrm{trop}}^*\alpha$ is, and
that is the next point.

\begin{lemma}
\label{lem:ch13-surjective-pullback-injective}
Let $Q\subset N'_{\mathbf R}$ and $P\subset N_{\mathbf R}$ be closed
polyhedral subsets, and let $F:N'_{\mathbf R}\to N_{\mathbf R}$ be an
integer-affine map with $F(Q)=P$.  Then
\[
 F^*:A^{p,q}(P)\longrightarrow A^{p,q}(Q)
\]
is injective for all $p,q$.
\end{lemma}

\begin{proof}
Choose locally finite polyhedral decompositions $\mathscr D$ of $Q$ and
$\mathscr E_0$ of $P$.  The images $F(D)$, $D\in\mathscr D$, are polyhedra
contained in $P$, and by \Cref{prop:tropical-common-refinement} there is a
polyhedral decomposition $\mathscr E$ of $P$ refining $\mathscr E_0$ such
that each $F(D)$ is a union of cells of $\mathscr E$.  (Locally finitely many
cells are involved near any given point, so the construction is legitimate
in the locally finite setting.)

Let $\alpha\in A^{p,q}(P)$ with $F^*\alpha=0$; we must show that $\alpha$
lies in the kernel \eqref{eq:ch10-form-kernel} computed for $\mathscr E_0$.
By the subdivision invariance recorded after
\Cref{def:ch10-forms-on-support} it is enough to prove it for
$\mathscr E$.

Fix a cell $\sigma\in\mathscr E$ and a point $x\in\RelInt(\sigma)$.  Since
$F(Q)=P$, there is $D\in\mathscr D$ with $x\in F(D)$; as $F(D)$ is a union
of cells of $\mathscr E$ and contains the point $x$ of $\RelInt(\sigma)$, it
contains $\sigma$.  Hence the affine span of $\sigma$ is contained in the
affine span of $F(D)$, and therefore
\[
 N_{\sigma,\mathbf R}
 \subset\ \text{direction space of }F(D)
 \ =\ DF\bigl(N_{D,\mathbf R}\bigr),
\]
where $DF$ is the linear part of $F$; the last equality holds because $F$ is
affine and $D$ has $N_{D,\mathbf R}$ as direction space.

Now let $v_1,\ldots,v_p,w_1,\ldots,w_q\in N_{\sigma,\mathbf R}$ and choose
$\widetilde v_i,\widetilde w_j\in N_{D,\mathbf R}$ with
$DF(\widetilde v_i)=v_i$ and $DF(\widetilde w_j)=w_j$.  Pick $y\in D$ with
$F(y)=x$.  Evaluating the ambient representative of $F^*\alpha$ at $y$ on
these tangent vectors of $D$ gives
\[
 (F^*\alpha)_y(\widetilde v_1,\ldots;\widetilde w_1,\ldots)
 =\alpha_{x}(v_1,\ldots;w_1,\ldots),
\]
and the left-hand side vanishes because $F^*\alpha=0$ in $A^{p,q}(Q)$, that
is, because its ambient representative lies in the kernel
\eqref{eq:ch10-form-kernel} for $\mathscr D$.  Therefore
$\alpha_x(v_1,\ldots;w_1,\ldots)=0$ for every $x\in\RelInt(\sigma)$, and by
continuity of the coefficients for every $x\in\sigma$.  As $\sigma$ was
arbitrary, $\alpha=0$ in $A^{p,q}(P)$.
\end{proof}

\begin{warning}{Surjectivity is essential, injectivity is irrelevant}
\Cref{lem:ch13-surjective-pullback-injective} says nothing about
injectivity of $F$ and everything about surjectivity onto $P$.  A projection
$\mathbf R^2\to\mathbf R$ restricted to a square is far from injective, yet
pullback of forms along it is injective, because every tangent direction
downstairs is hit.  Conversely, if $F(Q)$ is a proper polyhedral subset of
$P$, then every $\alpha$ supported away from $F(Q)$ dies, and the kernel is
enormous.  This is why \Cref{def:ch13-chart-category} insists that $P$
contain $\Trop_f(V)$ but does not insist that it equal it: the freedom is
convenient, and the nullity criterion below removes it.
\end{warning}

\begin{proposition}
\label{prop:ch13-nullity}
Let $V$ be a \emph{compact} analytic domain of $X$, let $(f,P)$ be a
tropical chart datum on $V$, and let $\alpha\in A^{p,q}(P)$.  Then
\[
 f_{\mathrm{trop}}^*\alpha=0\ \text{ in }A^{p,q}_{\mathrm{pf}}(V)
 \iff
 \alpha|_{\Trop_f(V)}=0\ \text{ in }A^{p,q}\bigl(\Trop_f(V)\bigr).
\]
Consequently $\alpha\mapsto f_{\mathrm{trop}}^*\alpha$ induces an injection
\[
 A^{p,q}\bigl(\Trop_f(V)\bigr)\lhook\joinrel\longrightarrow
 A^{p,q}_{\mathrm{pf}}(V),
\]
whose image consists of the classes representable in the chart $f$.
\end{proposition}

\begin{proof}
Since $V$ is compact, $\Trop_f(V)$ is a compact polyhedral subset of
$T_{\mathrm{trop}}$ by \Cref{thm:analytic-tropical-polyhedrality}, and
$\bigl(f,\Trop_f(V)\bigr)$ is an object of $\mathscr T(V)$ dominating
$(f,P)$ through the identity.  Replacing $P$ by $\Trop_f(V)$ changes neither
side of the asserted equivalence, so we may and do assume
$P=\Trop_f(V)$.

If $\alpha=0$ in $A^{p,q}(P)$ then its class is zero, which is the trivial
implication.  Conversely, suppose $f_{\mathrm{trop}}^*\alpha=0$.  By
filteredness (\Cref{lem:ch13-charts-filtered}) there is an object $(g,Q)$ of
$\mathscr T(V)$ dominating $(f,P)$ through an integer-affine $F$, with
$F^*\alpha=0$ in $A^{p,q}(Q)$.  Replacing $Q$ by $\Trop_g(V)$, which again
only shrinks, we may assume $Q=\Trop_g(V)$.  Then
\[
 F(Q)=F\bigl(g_{\mathrm{trop}}(V)\bigr)=f_{\mathrm{trop}}(V)=P,
\]
so $F|_Q$ is a surjection of compact polyhedral sets and
\Cref{lem:ch13-surjective-pullback-injective} gives $\alpha=0$.
\end{proof}

\begin{warning}{The presheaf $A^{p,q}_{\mathrm{pf}}$ is not separated}
Compactness in \Cref{prop:ch13-nullity} cannot be dropped, and the failure
is not a technicality about closures.  Chambert-Loir and Ducros exhibit a
$G$-tropical space whose underlying set is $\mathbf R$ with the
\emph{discrete} topology, tropicalized by the identity chart
\cite[Remark~3.4.3]{CLD25}.  For $\alpha=\dprime x$ one has
$\alpha|_{\{t\}}=0$ for every singleton domain $\{t\}$, since a point has no
tangent vectors, so the class of $\alpha$ restricts to zero on each member of
a $G$-cover; but $\alpha$ itself is not zero in the colimit, because every
chart dominating the identity does so through a \emph{surjective} affine map
and \Cref{lem:ch13-surjective-pullback-injective} applies.  Hence
$A^{1,0}_{\mathrm{pf}}$ is not a separated presheaf, and sheafification in
\Cref{thm:ch13-sheaf} is not cosmetic.  The moral for Berkovich spaces is
that a chart may fail to see any tangent direction over a piece of $X$
without the corresponding form being zero on that piece; only compactness of
the domain forces the tropical image to be large enough to detect it.
\end{warning}

\subsection{The sheaves \texorpdfstring{$A^{p,q}$}{Apq}}
\label{sec:ch13-sheaves}

\begin{definition}[CLD form]
\label{def:ch13-cld-form}
For $p,q\geq0$, a \emph{smooth form of bidegree $(p,q)$} on an open set
$U\subset X$ is a section locally representable as
\[
                    \omega|_V=f_{\mathrm{trop}}^*\alpha,
\]
where $V\subset U$ is open, $f$ is a moment map on an analytic neighborhood
of $V$, and $\alpha$ is a smooth superform of bidegree $(p,q)$ near
$f_{\mathrm{trop}}(V)$.  Representatives are identified when their pullbacks
agree after passage to a common product moment map.  The resulting sheaf is
denoted $A_X^{p,q}$; compactly supported sections are denoted
$A_{X,c}^{p,q}$.
\end{definition}

Unravelled, \Cref{def:ch13-cld-form} says that $A^{p,q}_X$ is the sheaf on
the topological space $X$ associated to the presheaf
$U\mapsto A^{p,q}_{\mathrm{pf}}(U)$ of \Cref{def:ch13-presheaf}.  The
following theorem says that the sheafification does not destroy the chart
description, and identifies the stalks.

\begin{theorem}
\label{thm:ch13-sheaf}
Let $X$ be a topologically separated $K$-analytic space and $p,q\ge0$.
\begin{enumerate}
\item\label{itm:ch13-sheaf:stalk} For $x\in X$ the stalk is
  \[
   A^{p,q}_{X,x}
   =\varinjlim_{(V,f,P)}A^{p,q}(P),
  \]
  the filtered colimit over tropical chart data on analytic neighbourhoods
  $V$ of $x$.
\item\label{itm:ch13-sheaf:sections} For $U\subset X$ open,
  $A^{p,q}_X(U)$ is the set of families $(\omega_x)_{x\in U}$ with
  $\omega_x\in A^{p,q}_{X,x}$ which are locally represented by a single
  chart datum; the assignment $U\mapsto A^{p,q}_X(U)$ is a sheaf of
  $\mathbf R$-vector spaces, and
  $A^{\bullet,\bullet}_X=\bigoplus_{p,q}A^{p,q}_X$ is a sheaf of bigraded
  $\mathbf R$-algebras.
\item\label{itm:ch13-sheaf:faithful} If $V$ is a compact analytic domain of
  $X$ and $(f,P)$ is a chart datum on $V$, the map
  $\alpha\mapsto f_{\mathrm{trop}}^*\alpha$ induces an injection
  $A^{p,q}(\Trop_f(V))\hookrightarrow A^{p,q}_{X,G}(V)$.
\item\label{itm:ch13-sheaf:GvsO} The natural map
  $A^{p,q}_X(U)\to A^{p,q}_{X,G}(U)$ is injective for every open
  $U\subset X$.
\end{enumerate}
\end{theorem}

\begin{proof}
\cref{itm:ch13-sheaf:stalk} Sheafification does not change stalks, and the
stalk of the presheaf at $x$ is the colimit of $A^{p,q}_{\mathrm{pf}}(V)$
over neighbourhoods $V$ of $x$, hence the displayed double colimit.  It is
filtered: neighbourhoods of $x$ are directed under inclusion, and the chart
data on each are directed by \Cref{lem:ch13-charts-filtered}, the
compactness hypothesis there being met because every point of a good space
has a basis of compact analytic neighbourhoods and, in general, one may
compute the stalk on a $G$-basis of compact analytic domains.

\cref{itm:ch13-sheaf:sections} This is the standard description of the
sheafification as the sheaf of locally coherent families of germs.  That the
sections form an $\mathbf R$-vector space is clear.  For the algebra
structure, given $\omega,\eta$ represented near $x$ in charts $(f,P)$,
$(g,Q)$, pass to the product chart of
\Cref{lem:ch13-charts-filtered}\,\cref{itm:ch13-cf:prod} and take the wedge
product of the two pullbacks there; the result is independent of the choices
because pullback along an integer-affine map is an algebra homomorphism
(\Cref{prop:ch09-pullback-functorial}) and because the colimit is
filtered.

\cref{itm:ch13-sheaf:faithful} By \Cref{prop:ch13-nullity} the map
$A^{p,q}(\Trop_f(V))\to A^{p,q}_{\mathrm{pf}}(V)$ is injective, so it
suffices to prove that
$A^{p,q}_{\mathrm{pf}}(V)\to A^{p,q}_{X,G}(V)$ is injective on the image,
that is, that a chart representative which is $G$-locally zero is zero.  Let
$f_{\mathrm{trop}}^*\alpha$ restrict to zero on each member of a $G$-cover of
the compact domain $V$; refining, we may take the cover to consist of
finitely many compact analytic domains $V_1,\ldots,V_m$.  By
\Cref{prop:ch13-nullity} applied on $V_i$, we get
$\alpha|_{\Trop_f(V_i)}=0$ for every $i$.  Now
$\Trop_f(V)=\bigcup_i\Trop_f(V_i)$, a finite union of compact polyhedral
subsets; choose a polyhedral decomposition of $\Trop_f(V)$ in which each
$\Trop_f(V_i)$ is a union of cells (\Cref{prop:tropical-common-refinement}).
For a cell $\sigma$ of that decomposition, pick $i$ with
$\sigma\subset\Trop_f(V_i)$; the vanishing condition
\eqref{eq:ch10-form-kernel} for $\alpha$ along $\sigma$ is then part of
$\alpha|_{\Trop_f(V_i)}=0$.  Since the kernel is defined by a pointwise
cellwise condition and is unchanged by subdivision, $\alpha|_{\Trop_f(V)}=0$,
whence $f_{\mathrm{trop}}^*\alpha=0$.

\cref{itm:ch13-sheaf:GvsO} Injectivity may be checked on germs, and even on
compact analytic neighbourhoods.  Let $\omega\in A^{p,q}_X(U)$ have zero
image.  Near a point $x$ write $\omega=f_{\mathrm{trop}}^*\alpha$ on an open
$W\ni x$, and choose a compact analytic neighbourhood $V\subset W$ of $x$;
then $f_{\mathrm{trop}}^*(\alpha)|_V$ is zero as a $G$-form, so
$\alpha|_{\Trop_f(V)}=0$ by \cref{itm:ch13-sheaf:faithful}, so
$\omega|_{\Int(V)}=0$.  As $x$ was arbitrary, $\omega=0$.
\end{proof}

\begin{corollary}
\label{cor:ch13-recognition}
Let $V$ be a compact analytic domain of $X$ and let $f:V\to T$ be a moment
map.  A $G$-form $\omega\in A^{p,q}_{X,G}(V)$ which is tropicalized by $f$
is of the form $f_{\mathrm{trop}}^*\alpha$ for a unique
$\alpha\in A^{p,q}(\Trop_f(V))$.  In particular, if $f_{\mathrm{trop}}$ is
injective on $V$ --- for instance if $V$ is a local skeleton on which the
chart is an isomorphism onto its image --- then $\omega$ determines and is
determined by an ordinary superform on the polyhedron $\Trop_f(V)$.
\end{corollary}

\begin{proof}
Existence is the hypothesis and uniqueness is
\Cref{thm:ch13-sheaf}\,\cref{itm:ch13-sheaf:faithful}.
\end{proof}

There is an analogous sheaf $A_{X,G}^{p,q}$ in which one uses analytic
domains and $G$-covers.  Ordinary forms are the objects used against the
usual topology of $X$; $G$-forms are indispensable for compact analytic
domains and boundary formulae.  On a good space, local questions can usually
be reduced to either language, but the two sheaves should not be identified
without stating the required comparison.

The precise relation is the one just proved: $A^{p,q}_X\to A^{p,q}_{X,G}$ is
injective, so an ordinary form is a $G$-form and nothing is lost by working
$G$-locally; but the map is not surjective, because a $G$-cover of an open
set by compact analytic domains need not refine to an open cover, and a
$G$-form assembled from charts on such domains need not be representable on
any ordinary neighbourhood of a boundary point.  This is exactly the
situation on a closed annulus viewed inside the affine line: the two Shilov
points have no ordinary neighbourhood inside the annulus, so a $G$-form built
from the annulus chart is not, without further argument, an ordinary form of
the ambient space.  The distinction is invisible for boundaryless spaces and
unavoidable for boundary formulae, which is why
\Cref{sec:ch13-stokes} is stated for $G$-forms.

\begin{proposition}\label{prop:cld-local-normal-form}
Every form $\omega\in A_X^{p,q}(U)$ admits an open cover $(V_i)$ and
representations
\[
                  \omega|_{V_i}=f_{i,\mathrm{trop}}^*\alpha_i.
\]
For a smooth function $u$ this says, more concretely, that near each point
there are invertible analytic functions $f_1,\ldots,f_r$ and a smooth real
function $\varphi$ such that
\[
 u=\varphi(-\log|f_1|,\ldots,-\log|f_r|).
\]
\end{proposition}
\begin{proof}
A section of $A_X^{p,q}$ is, by construction, locally an equivalence class
of a pair consisting of a moment map and a superform.  Choose one pair at
each point and shrink its analytic domain to an ordinary open neighborhood;
these neighborhoods give the asserted cover and representations.  In
bidegree $(0,0)$ the superform is an ordinary smooth function $\varphi$ on
the tropical target, which gives the displayed formula.

For the $G$-version, make the same choice on analytic domains.  If two such
pairs represent the same germ, their pullbacks coincide on a product moment
map by the equivalence relation, so the construction is independent of the
choice.  The comparison with the reduction is stated separately in
\Cref{prop:ch13-pullback-properties}\,
\cref{itm:ch13-pb:red}; its converse direction uses the almost-good lifting
theorem and is not part of the local-normal-form assertion.
\end{proof}

\begin{example}[An annulus]
\label{ex:ch13-annulus-radial}
Let $A=\{r<|T|<R\}$ and put $t=-\log|T|$.  In the chart given by $T$,
\[
 u=\varphi(t),\qquad
 \dprime u=\varphi'(t)\dprime t,\qquad
 \ddprime u=\varphi'(t)\ddprime t,
\]
and
\[
                    \dprimeddprime u
                     =\varphi''(t)\dprime t\wedge\ddprime t.
\]
Thus a smooth radial function is psh precisely when $\varphi$ is convex.
The formula also fixes the sign used in
\Cref{chap:cld-poincare-lelong,chap:cld-bedford-taylor}.
\end{example}

\subsection{Forms on the Berkovich affine line}
\label{sec:ch13-line}

Before any general theory is used, it is worth writing out every sheaf
$A^{p,q}$ on the simplest Berkovich space.  The computation is short and it
fixes the intuition for the rest of the chapter: forms live on the tree, not
on the disc.

\begin{example}[The sheaves $A^{p,q}$ on $\mathbf A^{1,\An}$]
\label{ex:ch13-affine-line}
Let $K$ be algebraically closed and nontrivially valued and let
$X=\mathbf A^{1,\An}_K$, of pure dimension $1$, with coordinate $T$.  We use
the classification of \Cref{thm:berk-a1-classification} and the tree
structure of \Cref{prop:berk-a1-tree}.

\emph{Vanishing.} $A^{p,q}_X=0$ as soon as $\max(p,q)\ge2$.  Indeed a form
of bidegree $(p,q)$ vanishes at every point $x$ with $d_K(x)<\max(p,q)$, and
$d_K(x)\le\dim_xX=1$ everywhere; alternatively, in any chart $f$ the tropical
image has dimension $\le1$ by \Cref{thm:trop-dimension-formula}, so a
superform of bidegree $(p,q)$ with $\max(p,q)\ge2$ has zero class in
$A^{p,q}(\Trop_f(V))$ by \eqref{eq:ch10-form-kernel}.  So only
$(0,0),(1,0),(0,1),(1,1)$ occur.

\emph{Functions.} Let $u$ be a smooth function near a point $x$.  By
\Cref{prop:cld-local-normal-form} there are invertible analytic functions
$f_1,\ldots,f_r$ near $x$ and a smooth $\varphi$ on $\mathbf R^r$ with
$u=\varphi(-\log|f_1|,\ldots,-\log|f_r|)$.  Each generator
$-\log|f_i|=v(f_i)$ is, on the skeleton of an affinoid domain on which $f_i$
is invertible, piecewise affine with integer slopes, by the slope formula
\Cref{thm:ch08-slope-formula}; so the local model of a smooth function on the
line is a smooth function of finitely many piecewise affine coordinates,
exactly as in the tropical picture.

\emph{Constancy on discs.} If $D=\{|T-a|\le\varrho\}$ is a closed disc
contained in the domain of definition of an invertible analytic function $f$,
then $|f|$ is constant on $D$.  Indeed $f$ has no zero on $D$, so writing
$f=c(1+h)$ with $c=f(a)\in K^\times$, the value at the centre, and
$\|h\|_D<1$ --- possible because
$K$ is algebraically closed, by Weierstrass division on the disc --- gives
$|f(y)|=|c|$ for all $y\in D$.  Consequently a smooth function is constant on
every closed disc contained in a single chart domain.  This is the
one-dimensional shadow of the skeletal comparison of
\Cref{sec:ch13-skeleton}, and it is why the sheaf $A^{0,0}$ is much smaller
than the sheaf of continuous functions: the function $x\mapsto|T(x)|$ is
continuous on $\mathbf A^{1,\An}$ but is \emph{not} smooth at the origin,
since it is not constant on any disc $\{|T|\le\varrho\}$ around it.  Away
from the origin it is smooth, being $\rme^{-t}$ in the chart $t=-\log|T|$.
The same example, in the $G$-topology, is
\cite[\S3.4.8]{CLD25}.

\emph{Bidegree $(1,1)$.} On the open subset
$\{0<|T|<\infty\}=\mathbf G_{m}^{\An}$ take the chart $f=T$, so
$t=-\log|T|$ is an affine coordinate on $\Trop_f(\mathbf G_m^{\An})
=\mathbf R$.  Then
\[
 A^{1,1}_X\big|_{\mathbf G_m^{\An}}
 =\{\rho(t)\,\dprime t\wedge\ddprime t\},
\]
with $\rho$ smooth; and the integral of a compactly supported such form is
$\int_{\mathbf R}\rho(t)\,\rd t$, the canonical calibration of the skeleton
$S(\mathbf G_m)$ being the constant $1$-vector
(\Cref{ex:trop-torus-explicit}).  A type-$1$ point contributes nothing:
$d_K=0$ there, and the tropical image of a small disc around it is a point.
\end{example}

The last computation deserves to be repeated at a point where the tree
branches, because that is where the difference between an analytic space and
a real interval becomes visible.

\begin{example}[A form supported near a type-$2$ point]
\label{ex:ch13-type-two}
Keep $K$ algebraically closed and nontrivially valued, and let
$X=\mathbf P^{1,\An}_K$.  Let $\eta$ be the Gauss point, that is, the
type-$2$ point corresponding to the closed unit disc.  Work on the open
subset
\[
 U=\mathbf P^{1,\An}\setminus\{0,1,\infty\}
\]
with the moment map $f=(T,\,T-1):U\to\mathbf G_m^2$, so that
\[
 f_{\mathrm{trop}}=(-\log|T|,\,-\log|T-1|).
\]
The image $\Trop_f(U)$ is the tropical line of
\Cref{ex:ch10-standard-line}: the union of the three rays
\[
 \mathbf R_{\ge0}u_1,\quad \mathbf R_{\ge0}u_2,\quad \mathbf R_{\ge0}u_3,
 \qquad
 u_1=(1,0),\ u_2=(0,1),\ u_3=(-1,-1),
\]
emanating from the origin, and $f_{\mathrm{trop}}(\eta)=(0,0)$.  Indeed if
$|T(x)|>1$ then $|T(x)-1|=|T(x)|$, giving the ray in direction
$(-1,-1)$; if $|T(x)|<1$ then $|T(x)-1|=1$, giving the ray $u_1$; and if
$|T(x)|=1$ then $|T(x)-1|\le1$, giving the ray $u_2$.  The three rays
correspond to the three ends of $U$, namely $\infty$, $0$ and $1$.

By \Cref{thm:trop-sigma-characterization} the local skeleton $\Sigma_f$ is the
set of $x\in U$ at which one of the initial forms $\widetilde{T(x)}$,
$\widetilde{(T-1)(x)}$ is transcendental over $\widetilde K$.  This is
exactly the tripod: the union of the three paths in the tree
$\mathbf P^{1,\An}$ joining $\eta$ to the ends $0$, $1$ and $\infty$, with
the three ends removed.  Indeed a point of the path from $\eta$ towards $0$
is $\eta_{0,\varrho}$ with $0<\varrho<1$, and then $\widetilde{T(x)}$ is
transcendental of degree $\varrho$; symmetrically for the other two paths,
using $T-1$ on the branch towards $1$.  Conversely, if $x=\eta_{a,\varrho}$
lies off the tripod, then $|a|=|a-1|=1$ and $\varrho<1$, so
$\widetilde{T(x)}=\widetilde a$ and
$\widetilde{(T-1)(x)}=\widetilde{a-1}$ both lie in $\widetilde K$, and
$d_x(f)=0$.  Type-$1$ and type-$4$ points have $d_K=0$ and are excluded as
well.

The map $f_{\mathrm{trop}}$ carries the tripod homeomorphically onto the
tropical line, with vertex $\eta\mapsto(0,0)$.  All calibration weights are
$1$: on the branch towards $\infty$ or $0$ the projection $q=\mathrm{pr}_1$
gives $q\circ f=T$, which realizes $U$ as an open subspace of
$\mathbf G_m^{\An}$ and is therefore of local degree $1$; on the branch
towards $1$ one uses $q=\mathrm{pr}_2$ and $T-1$ instead.  Hence
$\Pi_f$ is the tropical line with all weights $1$, which is balanced ---
this is \Cref{ex:ch10-standard-line} --- as
\Cref{thm:analytic-tropical-balancing} demands.

Now let $\chi$ be a smooth function on $\mathbf R^2$, supported in the
cube $\{\max_i|x_i|<\varepsilon\}$ with $\varepsilon<\tfrac12$, with
$\chi\equiv1$ near the origin, and set
\[
 \theta=f_{\mathrm{trop}}^*\bigl(\chi\,\ddprime x_1\bigr)
 \ \in A^{0,1}_{X}(U).
\]
Its support is the union of the three initial segments of the tripod of
lattice length $\varepsilon$, together with all the discs retracting onto
them.  So a form ``supported near $\eta$'' is not supported on a
neighbourhood of $\eta$ in any real chart: it is supported on a neighbourhood
of $\eta$ in $X$, whose tropical shadow is a \emph{star with three branches},
not an interval.

This support is a compact subset of $U$: on the branch in direction $u_1$ it
is $\{\rme^{-\varepsilon}\le|T|\le1,\ |T-1|=1\}$, and similarly for the other
two, each of which is a compact analytic domain of $\mathbf P^{1,\An}$ not
meeting $\{0,1,\infty\}$.  Hence $\theta$ extends by zero to an element of
$A^{0,1}_{X,c}(\mathbf P^{1,\An})$, and the integrals below are integrals
over the whole of $X$.

The Stokes computation is now the tripod computation of
\Cref{ex:ch11-tripod-stokes} verbatim.  Restricting to the ray $\sigma_j$
with lattice parameter $t$ gives
$\theta|_{\sigma_j}=\langle(\chi,0),u_j\rangle\,\ddprime t
=\chi(tu_j)u_{j,1}\,\ddprime t$, so
\[
 \int_{\sigma_j}\dprime\theta
 =\bigl[\chi u_{j,1}\bigr]_{t=0}^{t=\infty}
 =-\chi(0)\,u_{j,1},
\]
and summing over $j$,
\[
 \int_X\dprime\theta
 =-\chi(0)\sum_{j=1}^3u_{j,1}
 =-\chi(0)\bigl(1+0-1\bigr)=0 .
\]
The cancellation is the balancing relation $u_1+u_2+u_3=0$ at the vertex,
which here is the statement that the three ends of $U$ carry total slope
zero --- the degree-zero divisor formula for $T$ on $\mathbf P^1$, in the
guise of \Cref{thm:analytic-tropical-balancing}.  Had we used a punctured
space with an unbalanced star, Stokes would have failed, exactly as in
\Cref{ex:ch11-tripod-unbalanced}; but no such star arises from an analytic
space, and that is the content of interior harmony.
\end{example}

\section{Operations, locality, and pullback}
\label{sec:ch13-operations}

All algebraic operations are defined on a representative and then descended:
\[
 \omega\wedge\eta=f_{\mathrm{trop}}^*(\alpha\wedge\beta),\quad
 \dprime\omega=f_{\mathrm{trop}}^*(\dprime\alpha),\quad
 \ddprime\omega=f_{\mathrm{trop}}^*(\ddprime\alpha).
\]
If the representatives initially use different charts, one first replaces
them by their product chart.

For the wedge product and the involution $J$ this is unproblematic, because
those operations are $C^\infty$-multilinear and pullback along an affine map
commutes with them tautologically.  The differentials require one extra
verification, which is easy to overlook: the definition of a form on a
polyhedral support is a quotient by the cellwise kernel
\eqref{eq:ch10-form-kernel}, and it must be checked that this kernel is
stable under $\dprime$ and $\ddprime$.  If it were not, the differentials
would not even be defined on $A^{p,q}(P)$, let alone on $A^{p,q}_X$.

\begin{lemma}
\label{lem:ch13-kernel-differential}
Let $\mathscr C$ be a rational polyhedral complex in $N_{\mathbf R}$, let
$W\subset N_{\mathbf R}$ be open, and let $K^{p,q}=K^{p,q}(W,\mathscr C)$ be
the kernel \eqref{eq:ch10-form-kernel}.  Then
\[
 \dprime K^{p,q}\subset K^{p+1,q},
 \qquad
 \ddprime K^{p,q}\subset K^{p,q+1},
 \qquad
 JK^{p,q}=K^{q,p},
\]
and $K^{\bullet,\bullet}$ is a two-sided ideal of $A^{\bullet,\bullet}(W)$.
Consequently $\dprime$, $\ddprime$, $J$ and the wedge product descend to
$A^{\bullet,\bullet}(\Omega)$ for every open $\Omega\subset|\mathscr C|$.
\end{lemma}

\begin{proof}
The ideal property and the statement about $J$ are immediate from the
pointwise definition: the condition in \eqref{eq:ch10-form-kernel} is the
vanishing of all evaluations on cell tangent vectors, which is preserved by
multiplying by an arbitrary superform and by exchanging the two exterior
factors.

For $\dprime$, fix a cell $\sigma\in\mathscr C$ of dimension $d$ and a point
$x_0\in\Omega\cap\sigma$.  Choose affine coordinates $x_1,\ldots,x_m$ on
$N_{\mathbf R}$ such that $x_1,\ldots,x_d$ restrict to coordinates on the
affine span of $\sigma$ and $x_{d+1},\ldots,x_m$ vanish on that span; this is
possible because $N_{\sigma,\mathbf R}$ is a direct summand of
$N_{\mathbf R}$ as a real vector space.  Write
$\alpha=\sum_{I,J}\alpha_{IJ}\,\dprime x_I\wedge\ddprime x_J$ with $I,J$
increasing multi-indices in $\{1,\ldots,m\}$.  The condition
$\alpha\in K^{p,q}$ says exactly that
\begin{equation}
\label{eq:ch13-kernel-coordinates}
 \alpha_{IJ}\equiv0\ \text{ on }\Omega\cap\sigma
 \qquad\text{whenever }I\cup J\subset\{1,\ldots,d\},
\end{equation}
since evaluating $\alpha$ on tangent vectors of $\sigma$ picks out precisely
the coefficients whose indices lie in $\{1,\ldots,d\}$.

Now
$\dprime\alpha
=\sum_{I,J}\sum_{i=1}^m\partial_i\alpha_{IJ}\,
\dprime x_i\wedge\dprime x_I\wedge\ddprime x_J$, and its coefficient
attached to a pair of multi-indices contained in $\{1,\ldots,d\}$ is, up to
sign, a partial derivative $\partial_i\alpha_{IJ}$ with $i\le d$ and
$I\cup J\subset\{1,\ldots,d\}$.  For such indices, $\alpha_{IJ}$ vanishes
identically on the relatively open subset $\Omega\cap\sigma$ of the affine
span of $\sigma$, and $\partial_i$ with $i\le d$ differentiates along that
span; hence $\partial_i\alpha_{IJ}=0$ on $\Omega\cap\RelInt(\sigma)$, and by
continuity on $\Omega\cap\sigma$.  Thus $\dprime\alpha$ satisfies
\eqref{eq:ch13-kernel-coordinates} for $\sigma$.  As $\sigma$ was arbitrary,
$\dprime\alpha\in K^{p+1,q}$.  The assertion for $\ddprime$ follows from
$\ddprime=J\dprime J$ (\Cref{lem:ch09-J-conjugation}) together with
$JK^{p,q}=K^{q,p}$.
\end{proof}

\begin{warning}{Why a differential operator needs the tangent lattice}
\Cref{lem:ch13-kernel-differential} is the reason a form on a polyhedral
support may be differentiated at all, and it also shows the exact price:
only derivatives \emph{along} cells survive.  A superform whose coefficients
vary rapidly in a normal direction represents the zero form, and so does its
differential.  In particular $\dprime$ on $A^{p,q}(\Omega)$ is not computed
by an ambient gradient; it is computed by the gradient tangential to each
cell, and the two disagree at every point of a cell of positive
codimension.  On a Berkovich space this is why $\dprimeddprime u$ of a
smooth function is supported on the local skeleta and not on the whole
space.
\end{warning}

\begin{proposition}\label{prop:cld-intrinsic-calculus}
Wedge product, the involution $J$, $\dprime$, and $\ddprime$ are independent
of the chosen representatives.  They satisfy
\[
 (\dprime)^2=(\ddprime)^2=0,
 \qquad \dprime\ddprime+\ddprime\dprime=0,
\]
the graded Leibniz rules, and restriction to open subsets.  Hence
$A_X^{\bullet,\bullet}$ is a sheaf of bidifferential bigraded algebras.
\end{proposition}
\begin{proof}
Two representatives are equal precisely after pullback to a common
refinement.  Pullback of superforms by an affine map commutes with wedge,
$J$, $\dprime$, and $\ddprime$.  The asserted identities therefore follow
from their affine counterparts in \Cref{chap:superforms}.  Locality
and gluing follow from the definition of equality in refinements.
\end{proof}

\subsection{Locality}
\label{sec:ch13-locality}

The word ``local'' is used in three different senses in this subject, and the
proposition below separates them.  A sheaf-theoretic operator is local if it
commutes with restriction; a differential operator is local if it does not
increase supports; and an operator is $C^\infty$-linear if it is a
morphism of $A^{0,0}$-modules.  The differentials $\dprime$ and $\ddprime$
have the first two properties and not the third.

\begin{proposition}
\label{prop:ch13-locality}
Let $U\subset X$ be open and $\omega\in A^{p,q}_X(U)$.
\begin{enumerate}
\item\label{itm:ch13-loc:restrict} For every open $U'\subset U$ and every
  analytic domain $V\subset U$ one has
  $\dprime(\omega|_{U'})=(\dprime\omega)|_{U'}$ and
  $\dprime(\omega|_{V})=(\dprime\omega)|_{V}$ in $A^{p+1,q}_{X,G}(V)$, and
  likewise for $\ddprime$, $J$ and $\wedge$.  Thus $\dprime,\ddprime$ are
  morphisms of sheaves of $\mathbf R$-vector spaces.
\item\label{itm:ch13-loc:support} $\Supp(\dprime\omega)\subset\Supp\omega$
  and $\Supp(\ddprime\omega)\subset\Supp\omega$; more generally
  $\Supp(\omega\wedge\eta)\subset\Supp\omega\cap\Supp\eta$.
\item\label{itm:ch13-loc:leibniz} For $u\in A^{0,0}_X(U)$,
  \[
   \dprime(u\omega)=\dprime u\wedge\omega+u\,\dprime\omega,
  \]
  so $\dprime$ is not $A^{0,0}_X$-linear; it is a first-order differential
  operator, and its symbol is $u\mapsto\dprime u$.
\item\label{itm:ch13-loc:dim} $\Supp\omega\subset\{x\in U:d_K(x)\ge\max(p,q)\}$
  in the notation of \Cref{def:trop-residual-degree}; in particular
  $A^{p,q}_X=0$ as soon as $\dim X<\max(p,q)$.
\end{enumerate}
\end{proposition}

\begin{proof}
\cref{itm:ch13-loc:restrict} All four operations are defined chartwise, and a
chart datum on $U$ restricts to a chart datum on $U'$ or on $V$ by
restricting the moment map, with the same target polyhedron.  Since the
operations on $A^{p,q}(P)$ do not involve the domain, the two computations
agree.  Independence of the representative is
\Cref{prop:cld-intrinsic-calculus}.

\cref{itm:ch13-loc:support} If $\omega|_{U'}=0$ then
$(\dprime\omega)|_{U'}=\dprime(\omega|_{U'})=0$ by
\cref{itm:ch13-loc:restrict}; taking $U'$ to be the complement of
$\Supp\omega$ gives the inclusion.  The statement for $\wedge$ is the same
argument applied to the two factors separately.

\cref{itm:ch13-loc:leibniz} Represent $u$ and $\omega$ in a common chart, so
$u=f_{\mathrm{trop}}^*\varphi$ and $\omega=f_{\mathrm{trop}}^*\alpha$.  The
Leibniz rule of \Cref{prop:superform-differential-identities} in bidegree
$(0,0)$ gives $\dprime(\varphi\alpha)=\dprime\varphi\wedge\alpha
+\varphi\,\dprime\alpha$ on the target polyhedron; pulling back preserves it.
Failure of $A^{0,0}$-linearity is exhibited by any nonconstant $u$ and
$\omega=1$.

\cref{itm:ch13-loc:dim} Let $x\in U$ with $d_K(x)<\max(p,q)$ and represent
$\omega=f_{\mathrm{trop}}^*\alpha$ on a neighbourhood $V_0$ of $x$, with
$\xi=f_{\mathrm{trop}}(x)$.  The dimension chain
\eqref{eq:trop-dimension-chain} gives
$\dim_{\mathrm{trop},x}(f)\le d_x(f)\le d_K(x)<\max(p,q)$, so
$\dim_\xi f_{\mathrm{trop}}(V_0)<\max(p,q)$ after shrinking $V_0$, by
\Cref{def:trop-tropical-dimension} and
\Cref{thm:trop-local-germ}\,\cref{itm:trop-germ:germ}.  Since
$f_{\mathrm{trop}}$ is continuous, there is a compact analytic neighbourhood
$V\subset V_0$ of $x$ with $f_{\mathrm{trop}}(V)$ contained in a ball about
$\xi$ so small that
$\dim\Trop_f(V)=\dim_\xi f_{\mathrm{trop}}(V_0)<\max(p,q)$.  Every cell of
$\Trop_f(V)$ then has dimension $<\max(p,q)$, so its tangent space cannot
carry $\max(p,q)$ linearly independent vectors, and $\alpha$ lies in the
kernel \eqref{eq:ch10-form-kernel}; hence $\omega|_{V}=0$ and
$x\notin\Supp\omega$.  Since $d_K(x)\le\dim_xX$ by
\eqref{eq:trop-dimension-chain}, the last assertion follows.
\end{proof}

Statement \cref{itm:ch13-loc:dim} is the vanishing recorded in the remark
following \Cref{prop:cld-form-pullback}; the proof just given derives it from
the dimension chain of \Cref{chap:analytic-tropicalization} rather than
importing it, and it shows where the improvement from the tropical dimension
to $d_K$ comes from: one is allowed to shrink the chart domain, and the germ
theorem \Cref{thm:trop-local-germ} says the germ of the tropicalization then
stabilizes at the value predicted by the graded residue field.

\subsection{Pullback along an analytic morphism}
\label{sec:ch13-pullback}

\begin{proposition}\label{prop:cld-form-pullback}
Let $g:Y\to X$ be a morphism of topologically separated analytic spaces.
There is a functorial morphism
\[
                 g^*:A_X^{p,q}\longrightarrow g_*A_Y^{p,q}
\]
commuting with wedge, $J$, $\dprime$, and $\ddprime$.
\end{proposition}
\begin{proof}
If $\omega=f_{\mathrm{trop}}^*\alpha$, then on the inverse image of the
chart set put $g^*\omega=(f\circ g)_{\mathrm{trop}}^*\alpha$.  A common
refinement for two representatives pulls back to a common refinement on
$Y$, proving independence.  Functoriality and compatibility with the
operators are chartwise identities.
\end{proof}

\begin{remark}[Vanishing by dimension]
At a point $x$, a form of bidegree $(p,q)$ vanishes if the local tropical
dimension is smaller than $\max\{p,q\}$; see
\cite[Proposition~15.2.1]{CLD25}.  More generally, if an analytic subspace
$i:Z\hookrightarrow X$ has local dimension smaller than $\max\{p,q\}$,
then $i^*\omega=0$ for every $\omega\in A_X^{p,q}$.  This is a statement
about restriction to $Z$: the ambient support of $\omega$ may meet, or even
contain, $Z$.  In particular, dimension vanishing alone gives no
removability or uniqueness theorem for currents across $Z$.
\end{remark}

We record the further formal properties of $g^*$, since they are the ones
used in \Cref{chap:cld-currents,chap:cld-metrized-bundles} and since one of
them --- compatibility with extension of scalars --- is what makes the theory
insensitive to enlarging $K$.

\begin{proposition}
\label{prop:ch13-pullback-properties}
Let $g:Y\to X$ and $h:Z\to Y$ be morphisms of topologically separated
$K$-analytic spaces.
\begin{enumerate}
\item\label{itm:ch13-pb:funct} $(g\circ h)^*=h^*\circ g^*$ and
  $\mathrm{id}^*=\mathrm{id}$; in particular $A^{p,q}_{(-)}$ is a
  contravariant functor to sheaves.
\item\label{itm:ch13-pb:domain} If $g$ is the inclusion of an analytic
  domain, $g^*$ is the restriction map $A^{p,q}_{X,G}\to A^{p,q}_{Y,G}$.
\item\label{itm:ch13-pb:supp} $\Supp(g^*\omega)\subset g^{-1}(\Supp\omega)$.
\item\label{itm:ch13-pb:base} If $L/K$ is a complete valued extension and
  $p:X_L\to X$ the base-change morphism, then
  $p^*:A^{p,q}_X\to p_*A^{p,q}_{X_L}$ is injective, and every chart datum on
  $X$ pulls back to a chart datum on $X_L$ with the same target polyhedron.
\item\label{itm:ch13-pb:red} The canonical map
  $A^{p,q}_X\to A^{p,q}_{X_{\mathrm{red}}}$ induced by the closed immersion
  $X_{\mathrm{red}}\hookrightarrow X$ is an isomorphism on a good space.
\end{enumerate}
\end{proposition}

\begin{proof}
\cref{itm:ch13-pb:funct} On representatives, $(f\circ g\circ h)_{\mathrm{trop}}
=((f\circ g)\circ h)_{\mathrm{trop}}$, since composing a moment map with a
morphism of analytic spaces is associative and tropicalization is functorial
(\Cref{lem:trop-moment-basic}\,\cref{itm:trop-moment:funct}).  The identity
morphism carries a chart to itself.

\cref{itm:ch13-pb:domain} If $f$ is a moment map on an analytic domain $W$ of
$X$, then $f\circ g$ is its restriction to $W\cap Y$, with the same target
and the same tropical map; the induced map on colimits is the restriction map
of \Cref{def:ch13-presheaf}.

\cref{itm:ch13-pb:supp} If $\omega$ vanishes on an open $U$, then
$g^*\omega$ vanishes on $g^{-1}(U)$ by \cref{itm:ch13-pb:funct} and
\cref{itm:ch13-pb:domain}.

\cref{itm:ch13-pb:base} Tropicalization commutes with base change
(\Cref{lem:trop-moment-basic}\,\cref{itm:trop-moment:base}), and
$\Trop_{f_L}(V_L)=\Trop_f(V)$ for a compact analytic domain, since $p$ is
surjective and $|f_i(x_L)|=|f_i(p(x_L))|$.  Thus a chart datum pulls back to
one with the same polyhedron, and if $p^*\omega=0$ then, computing on a
compact analytic domain and applying
\Cref{thm:ch13-sheaf}\,\cref{itm:ch13-sheaf:faithful} over $L$, the
representative $\alpha$ vanishes on $\Trop_{f_L}(V_L)=\Trop_f(V)$, whence
$\omega=0$.

\cref{itm:ch13-pb:red} Restriction of an ordinary moment chart from $X$
to $X_{\mathrm{red}}$ preserves the absolute values of all coordinates and
hence its tropical image and superform representative.  Conversely, $X$
is good and therefore almost good.  The ordinary-open assertion of
\cite[Proposition~15.1.3(b) and the form variant immediately following
it]{CLD25} lifts each finite tuple of invertible-function germs on
$X_{\mathrm{red}}$ to unit germs on ordinary open neighbourhoods in $X$.
After shrinking and taking product moment maps on overlaps, these lifted
charts and their refinements form a cofinal system.  The two colimits
therefore agree.  The Gerritzen--Grauert argument in part~(a) concerns
$G$-domains and gives the analogous $G$-form comparison; it does not
replace the ordinary-open lifting used here.
\end{proof}

\subsection{What is, and what is not, a pullback}
\label{sec:ch13-what-is-pullback}

The symbol $f_{\mathrm{trop}}^*$ and the symbol $g^*$ of
\Cref{prop:cld-form-pullback} are both called pullbacks, and they are not the
same kind of object.  Confusing them is the commonest source of false
statements in this theory, so we separate them explicitly.

\paragraph{$g^*$ is a genuine sheaf-theoretic pullback.}
For a morphism $g:Y\to X$ of analytic spaces,
\Cref{prop:cld-form-pullback} produces a map of sheaves
$A^{p,q}_X\to g_*A^{p,q}_Y$ on $X$, equivalently a map
$g^{-1}A^{p,q}_X\to A^{p,q}_Y$ on $Y$ by the adjunction between $g^{-1}$ and
$g_*$.  It is functorial, compatible with all operations, and local.  In this
respect the sheaves $A^{p,q}$ behave exactly like sheaves of differential
forms on manifolds.  Nothing more is claimed, and in particular:

\begin{warning}{Three things $g^*$ is not}
\emph{It is not a map of $\mathcal O$-modules.}  The sheaves $A^{p,q}_X$ are
sheaves of $\mathbf R$-vector spaces and of $A^{0,0}_X$-modules only;
$\mathcal O_X$ does not act on them, and there is no coherent sheaf whose
inverse image they compute.  Consequently there is no ``$g^*$ of a module''
in the sense of $g^{-1}(-)\otimes_{g^{-1}\mathcal O_X}\mathcal O_Y$, and no
base-change spectral sequence.

\emph{It is neither injective nor surjective in general.}  Injectivity fails
for any $g$ whose image misses part of the support: if $Z\subset X$ is a
closed analytic subspace with $\dim Z<\max(p,q)$ and $i:Z\to X$ is the
inclusion, then $i^*\omega=0$ for every $\omega\in A^{p,q}_X$ by
\Cref{prop:ch13-locality}\,\cref{itm:ch13-loc:dim} applied on $Z$, no matter
how large $\Supp\omega$ is.  Surjectivity fails already for a closed
immersion of a point, and for the normalization of a nodal curve.

\emph{There is no pushforward of forms.}  A map $g_*:A^{p,q}_Y\to
A^{p,q}_X$ does not exist, not even for finite $g$: a form on $Y$ has no
canonical trace to $X$ because the chart data on $Y$ have no reason to come
from $X$.  What does exist is a pushforward of \emph{measures} and of
\emph{currents}, and only under properness hypotheses; see
\Cref{prop:ch13-finite-integral} and \Cref{chap:cld-currents}.
\end{warning}

\paragraph{$f_{\mathrm{trop}}^*$ is not a pullback at all.}
There is no map of ringed spaces, and no differentiable map, whose pullback
$f_{\mathrm{trop}}^*$ could be.  The map $f_{\mathrm{trop}}:V\to
T_{\mathrm{trop}}$ is merely continuous, its source is not a manifold, and it
has no differential.  What $f_{\mathrm{trop}}^*\alpha$ denotes is the class
of $\alpha$ in the filtered colimit of \Cref{def:ch13-presheaf}; the notation
is chosen because the resulting assignment is functorial in the chart,
$(F\circ f)_{\mathrm{trop}}^*\alpha=f_{\mathrm{trop}}^*(F^*\alpha)$, which is
the only property ever used.  \Cref{prop:ch13-nullity} measures exactly how
far it is from being injective, and the answer --- that the kernel is
$\{\alpha:\alpha|_{\Trop_f(V)}=0\}$ and nothing more --- is a theorem, not a
definition.

\begin{warning}{Two failures of the pullback intuition}
\emph{The image of a chart is not a submanifold.}  A tropical chart is not a
homeomorphism onto its image and its fibres are large: over an interior point
of $\Trop_f(V)$ sits an entire analytic fibre.  A form is therefore not
determined by its values ``at points of $X$'' in any naive sense; it is
determined by its restriction to the local skeleta, which is the content of
\Cref{prop:ch13-skeletal-faithful}.

\emph{Composition of charts is not composition of maps.}  If $(V',\psi)$
refines $(V,\varphi)$ through $F$, then $F$ relates the two \emph{targets},
not the two domains, and the refinement relation is directed and not
symmetric (\Cref{lem:trop-refinement-calculus}).  Two charts are compared by
passing up to their product, never down to a common coarsening, and the
resulting weight comparisons are pushforward statements carrying the lattice
index \eqref{eq:appG-lattice-index}.
\end{warning}

\section{The test-form topology}
\label{sec:ch13-topology}

Currents require more than the vector space $A_{X,c}^{p,q}$: they require its
locally convex topology.  We spell it out because weak convergence of
currents always refers to this choice.

Fix a compact set $C\subset X$ and finitely many tropical charts covering
$C$.  Refine their tropical images by finite polyhedral complexes and choose
affine coordinates on all cells.  For a chartwise representative $\alpha$,
let
\[
 \|\alpha\|_{C,r}
 =\max\bigl\{|\partial^I\alpha_{KL}(x)|:
      |I|\leq r,\ x\text{ in a relevant cell},\ K,L\bigr\}.
\]
Changing the finite refinement or affine coordinates gives an equivalent
family of seminorms.  The fixed-support, fixed-chart space is given the
$C^\infty$ topology generated by all $\|\cdot\|_{C,r}$.  Finally,
$A_{X,c}^{p,q}$ is given the final locally convex topology as the support and
finite chart data vary.

\begin{proposition}\label{prop:test-form-continuity-criterion}
A linear functional $T$ on $A_{X,c}^{p,q}$ is continuous if and only if its
restriction to every fixed compact-support, finite-chart space is continuous.
Equivalently, for each such datum there exist $C>0$, an integer $r$, and
finitely many compact chart domains $W_i$ such that
\[
                |\langle T,\omega\rangle|
               \leq C\max_i\|\omega\|_{W_i,r}.
\]
The integer $r$ may depend on both the compact set and the chart datum.
\end{proposition}
\begin{proof}
This is the defining universal property of the final topology.  A continuous
linear functional on a locally convex space controlled by a family of
seminorms is bounded by a finite maximum of those seminorms.  Conversely,
such bounds on every member of the final system imply continuity there and
hence on the final space.
\end{proof}

The dependence on the chart datum is real: even on a compact tropical space,
a current need not have one globally bounded distributional order.  This is
why later compactness assertions are stated in the weak topology rather than
in a fixed-order norm.

Two features of this topology deserve to be stated as a proposition, because
they are what make the operations of \Cref{sec:ch13-operations} usable
against currents.

\begin{proposition}
\label{prop:ch13-test-continuity}
Let $X$ be topologically separated.
\begin{enumerate}
\item\label{itm:ch13-tc:mult} For $u\in A^{0,0}_X(X)$ and
  $\gamma\in A^{r,s}_X(X)$, the maps $\omega\mapsto u\omega$ and
  $\omega\mapsto\omega\wedge\gamma$ are continuous
  $A^{p,q}_{X,c}\to A^{p,q}_{X,c}$, respectively
  $A^{p,q}_{X,c}\to A^{p+r,q+s}_{X,c}$.
\item\label{itm:ch13-tc:diff} $\dprime:A^{p,q}_{X,c}\to A^{p+1,q}_{X,c}$
  and $\ddprime$ are continuous, and increase by at most one the integer $r$
  occurring in \Cref{prop:test-form-continuity-criterion}.
\item\label{itm:ch13-tc:pullback} If $g:Y\to X$ is a morphism which is
  topologically proper, then $g^*:A^{p,q}_{X,c}\to A^{p,q}_{Y,c}$ is
  continuous.
\end{enumerate}
\end{proposition}

\begin{proof}
\cref{itm:ch13-tc:mult} On a fixed compact support and a fixed chart datum,
after enlarging the chart to a product tropicalizing $u$ or $\gamma$ as well,
multiplication is multiplication of superform coefficients on finitely many
cells; the Leibniz estimate
$\|ab\|_{C,r}\le c(r)\|a\|_{C,r}\|b\|_{C,r}$ gives continuity on each member
of the inductive system, hence on the limit.

\cref{itm:ch13-tc:diff} In the same fixed chart datum,
$\|\dprime\alpha\|_{C,r}\le\|\alpha\|_{C,r+1}$ by
\Cref{def:ch09-operators}, since the coefficients of $\dprime\alpha$ are
first partial derivatives of those of $\alpha$ in the chosen affine
coordinates.  The support does not grow, by
\Cref{prop:ch13-locality}\,\cref{itm:ch13-loc:support}.

\cref{itm:ch13-tc:pullback} Properness ensures that $g^{-1}$ of a compact set
is compact, so $g^*$ maps the fixed-support subspaces into fixed-support
subspaces.  In a chart, $g^*$ replaces the moment map $f$ by $f\circ g$
without changing the target polyhedron or the representative, so all the
seminorms $\|\cdot\|_{C,r}$ are literally unchanged.
\end{proof}

\begin{warning}{The order is not uniform}
\Cref{prop:test-form-continuity-criterion} permits the integer $r$ to depend
on the chart datum as well as on the compact set, and this cannot be
improved.  The reason is the same as in the real-analytic setting: the space
$A^{p,q}_{X,c}$ is an inductive limit over an uncountable directed family of
chart data, not merely over compact subsets, and the finite subfamilies grow
in two independent directions --- more coordinates, and finer cells.  A
functional bounded by $\|\cdot\|_{C,r}$ for a fixed $r$ on every chart datum
is much more special than a continuous one, and no compactness theorem below
produces such a bound.  The distinction is not academic: it is why weak
limits of currents are taken in the weak-$*$ topology and why the mass
bounds of \Cref{chap:cld-bedford-taylor} are stated as bounds on measures
rather than as bounds on distributional orders.
\end{warning}

\section{Positivity of analytic forms}
\label{sec:ch13-positivity}

Positivity of superforms was analysed in
\Cref{sec:ch09-positivity}: three cones
$A^{p,p}_{+,\mathrm s}\subset A^{p,p}_+\subset A^{p,p}_{+,\mathrm w}$,
mutually dual in the sense of \Cref{thm:ch09-positivity-duality}, coinciding
for $p\in\{0,1,n-1,n\}$ and distinct in between
(\Cref{ex:ch09-weak-not-positive,ex:ch09-positive-not-strong}).  Transporting
them to a Berkovich space requires only that each cone be preserved by
pullback along an integer-affine map, which is
\Cref{prop:ch09-positivity-inclusions} together with
\Cref{lem:ch09-pullback-eleven}.  The definition below is therefore
legitimate, and the verification is the same product-chart argument as
everywhere else in this chapter.

\begin{definition}
\label{def:ch13-positivity}
A symmetric $(p,p)$-form on $X$ is strongly positive, positive, or weakly
positive if it has that property in every tropical chart.  We write
\[
 A_{X,+,\mathrm{s}}^{p,p}\subset
 A_{X,+}^{p,p}\subset
 A_{X,+,\mathrm{w}}^{p,p}.
\]
\end{definition}

\begin{proposition}
\label{prop:ch13-positivity-basic}
The positivity cones are well defined, local, convex, closed under
multiplication by nonnegative smooth functions, and preserved by analytic
pullback.  A smooth function $u$ is psh in the smooth CLD sense exactly when
$\dprimeddprime u$ is a positive $(1,1)$-form.
\end{proposition}
\begin{proof}
Affine pullback preserves each positivity cone of superforms.  Common
refinements therefore prove chart independence and locality.  The remaining
claims follow from the corresponding linear-algebra statements for
superforms.  The last assertion is the definition of smooth psh.
\end{proof}

\begin{example}[Toric Hessian]
On an analytic torus, let $u=F\circ\operatorname{trop}$ with $F$ smooth.
Then
\[
 \dprimeddprime u
 =\sum_{i,j}F_{ij}\,
      \dprime x_i\wedge\ddprime x_j.
\]
It is positive exactly when the real Hessian $(F_{ij})$ is positive
semidefinite.  This is the smooth part of the toric real
Monge--Amp\`ere dictionary; nonsmooth convex $F$ require currents.
\end{example}

\begin{example}[Positivity on a curve is positivity of one coefficient]
\label{ex:ch13-positivity-curve}
Let $X$ be purely one-dimensional and let $\omega\in A^{1,1}_X(U)$ be
symmetric.  In a chart with tropical coordinate $t$ on a one-dimensional
cell, $\omega=\varrho(t)\,\dprime t\wedge\ddprime t$, and the three
positivity notions of \Cref{def:ch13-positivity} coincide and amount to
$\varrho\ge0$; this is the case $p=1=n$ of
\Cref{prop:ch09-positivity-inclusions}, in which the three cones agree.  On
the Berkovich affine line with the chart $t=-\log|T|$ this says that a smooth
function $u=\varphi(t)$ has $\dprimeddprime u\ge0$ exactly when $\varphi$ is
convex, recovering \Cref{ex:ch13-annulus-radial}, and the associated Radon
measure of \Cref{cor:ch13-measure} is $\varphi''(t)\,\rd t$ on the skeleton.
The atomic measures produced by piecewise affine $\varphi$ are not of this
form and are the first genuine currents of \Cref{chap:cld-currents}.
\end{example}

\section{Integration on the analytic space}
\label{sec:ch13-integration}

Assume now that $X$ is pure of dimension $n$.  Its tropicalizations carry
canonical calibrations, including analytic multiplicities and residue-field
degrees.

Before writing the definition we isolate the computation in a single chart,
because it is the formula used in every example and because it makes visible
which of the three inputs --- the polyhedron, the weights, the lattice
measure --- carries which piece of information.  Recall from
\Cref{def:trop-tropical-polyhedron} that a moment map $f$ on a space of level
$n$ is \emph{skeletally proper} if $f_{\mathrm{trop}}|_{\Sigma_f^{(n)}}$ is
topologically proper, which holds automatically for compact $X$, and that
$\Pi_f=f_{\mathrm{trop}}(\Sigma_f^{(n)})$ carries the pushforward of the
canonical calibration, with weights $m_D$ given by
\eqref{eq:trop-multiplicity-formula}.

\begin{proposition}
\label{prop:ch13-chart-integral}
Let $X$ be of level $n$ and let $f:X\to T$ be a skeletally proper moment map.
Let $\mathscr D$ be a polyhedral decomposition of $\Pi_f$ with weights
$(m_D)_{D\in\mathscr D_n}$ as in \Cref{thm:trop-multiplicity-formula}, and
let $\alpha$ be an $(n,n)$-superform on a polyhedral subset $P$ containing
$\Trop_f(X)$, whose support meets $\Pi_f$ in a compact set.  Then the
quantity
\begin{equation}
\label{eq:ch13-chart-integral}
 I_f(\alpha)
 \ =\ \int_{(\mathscr D,m)}\alpha
 \ =\ \sum_{D\in\mathscr D_n}m_D\int_D\alpha ,
\end{equation}
in which $\int_D\alpha=\int_Da\,\omega_D$ for $\alpha|_D=a\,\tau_n$ written
in lattice coordinates on $D$, as in \Cref{sec:ch11-top-degree},
depends only on the class $f_{\mathrm{trop}}^*\alpha\in A^{n,n}_{X,G}(X)$,
not on $P$, on $\alpha$, or on $\mathscr D$.  It is nonnegative if $\alpha$
is a positive $(n,n)$-form, and it vanishes if $\Supp\alpha$ does not meet
$\Pi_f$.
\end{proposition}

\begin{proof}
Independence of $\mathscr D$ is \Cref{prop:integral-subdivision-independent}
together with \Cref{lem:trop-calibration-subdivision}: a subdivision of
$\mathscr D$ inherits the same weights and the same lattice measures on
$n$-cells, by \Cref{lem:ch10-subdivision-lattices}.

Independence of $P$ and $\alpha$: suppose $f_{\mathrm{trop}}^*\alpha=0$.
Working $G$-locally on compact analytic domains and applying
\Cref{prop:ch13-nullity}, we get $\alpha|_{\Trop_f(V)}=0$ for every compact
analytic domain $V$, hence $\alpha|_{\Pi_f}=0$ because $\Pi_f$ is covered by
the $\Trop_f(V)$.  A cell $D\in\mathscr D_n$ has dimension $n$, so the
condition \eqref{eq:ch10-form-kernel} at a point of $D$ with tangent vectors
in $N_{D,\mathbf R}$ says precisely that the coefficient of $\alpha$ against
$\tau_n$ in lattice coordinates on $D$ vanishes; hence each summand of
\eqref{eq:ch13-chart-integral} is zero.  Linearity gives the general
statement.

Positivity: a positive $(n,n)$-form has nonnegative coefficient against
$\tau_n$ in every choice of coordinates (\Cref{prop:ch09-positivity-dictionary}),
and $m_D>0$ by \Cref{thm:trop-multiplicity-formula}.  The last assertion is
clear, since only cells of $\Pi_f$ contribute.
\end{proof}

\begin{remark}[What each factor in \eqref{eq:ch13-chart-integral} does]
\label{rem:ch13-three-factors}
The polyhedron $\Pi_f$ says \emph{where}; the lattice measure $\omega_D$
of \Cref{def:ch11-lattice-measure} says \emph{with what normalization}; and
the weight $m_D$ says \emph{how many times}.  Only the first is visible to
the naked eye, and it is the least canonical of the three: two charts have
different $\Pi_f$.  The lattice measure depends on the integral affine
structure and not on any Euclidean metric, which is why
\Cref{ex:ch11-nonunimodular} and \Cref{ex:ch11-euclidean-normal-trap} must be
kept in mind.  The weight is where the analytic geometry enters, through the
local degree of a finite flat projection
(\Cref{def:trop-canonical-calibration}); over a non-algebraically closed
field or a nonreduced space it is not a number of points, as
\Cref{ex:trop-multiplicity-two} and \Cref{ex:trop-length-factor} show.
\end{remark}

For $\omega\in A_{X,c}^{n,n}$, choose finitely many chart domains
$V_i$, a smooth partition of unity $(\rho_i)$, and representations
$\rho_i\omega=f_{i,\mathrm{trop}}^*\alpha_i$.  Define
\[
                       \int_X\omega
                  =\sum_i\int_{f_{i,\mathrm{trop}}(V_i)}
                            \alpha_i
\]
with the canonical calibrations.

Each summand is \eqref{eq:ch13-chart-integral} for the chart $f_i$; the sum
is finite because $\Supp\omega$ is compact and the partition is locally
finite (\Cref{thm:cld-partitions-unity}).  The theorem below is the assertion
that nothing in this recipe was a choice.

\begin{theorem}\label{thm:cld-analytic-integration}
The integral is independent of the cover, partition of unity, moment maps,
and subdivisions.  A locally integrable $(n,n)$-form defines a Radon measure.
Formation of this measure is compatible with complete extension of the
ground field.
\end{theorem}
\begin{proof}
Take two choices of covers, partitions, and moment maps.  Compact support
leaves finitely many chart pieces.  Intersect the two covers and replace the
two moment maps on every intersection by their product moment map.  After a
common rational polyhedral subdivision, all representatives live on one
finite family of calibrated complexes.  The tropical projection formula
\eqref{eq:analytic-tropical-pushforward} identifies the integral before and
after each projection from a product chart, while
\Cref{prop:integral-subdivision-independent} removes the subdivision.  The
two partitions disappear because their products form a common partition and
\(
                 \sum_{i,j}\rho_i\rho'_j=1.
\)
This proves independence.

We spell out the comparison step, since it is the one that uses the results
of \Cref{chap:analytic-tropicalization} rather than formal manipulation.
Let $V$ be a compact analytic domain carrying two moment maps $f:V\to T$ and
$g:V\to T'$, both tropicalizing $\omega$, and let $h=(f,g):V\to T\times T'$
be the product, with projections $p,p'$.  Then $h$ is again skeletally
proper and $\Sigma_f\cap V\subset\Sigma_h$ by
\Cref{lem:trop-refinement-calculus}\,\cref{itm:trop-ref:sigma}, so
$\Sigma_h^{(n)}=\Sigma_f^{(n)}\cap V$ by
\Cref{thm:trop-sigma-characterization}, both being the locus $d_x=n$.  By
\Cref{thm:trop-calibration-chart-independent} the canonical calibrations of
this common local skeleton computed with $f$ and with $h$ coincide.
Pushing forward along $f_{\mathrm{trop}}$, respectively
$h_{\mathrm{trop}}$, and using
$f_{\mathrm{trop}}=p_{\mathrm{trop}}\circ h_{\mathrm{trop}}$, we get
$p_{\mathrm{trop},*}\mu_{\Pi_h}=\mu_{\Pi_f}$, which is exactly
\eqref{eq:analytic-tropical-pushforward} with the lattice index appearing
once.  Since $\omega|_V=f_{\mathrm{trop}}^*\alpha
=h_{\mathrm{trop}}^*(p_{\mathrm{trop}}^*\alpha)$, the change-of-variables
identity for a piecewise immersion (\Cref{lem:ch11-cell-change-of-variables}
and \Cref{prop:ch11-pushforward-integral}) gives
\[
 \int_{(\mathscr D_h,m_h)}p_{\mathrm{trop}}^*\alpha
 =\int_{(\mathscr D_f,m_f)}\alpha ,
\]
that is, $I_h(p_{\mathrm{trop}}^*\alpha)=I_f(\alpha)$ in the notation of
\eqref{eq:ch13-chart-integral}.  The same computation with $g$ and $p'$ gives
$I_h((p')_{\mathrm{trop}}^*\beta)=I_g(\beta)$, and the two left-hand sides
agree because $p_{\mathrm{trop}}^*\alpha$ and $(p')_{\mathrm{trop}}^*\beta$
have the same class in the chart $h$, by
\Cref{prop:ch13-chart-integral}.  This is the required comparison.

In a chart, a locally integrable top form is a locally integrable density
with respect to lattice-normalized Lebesgue measure on finitely many
polyhedra.  Multiplication by the positive calibration weights and summation
therefore give a Radon measure.  The preceding independence shows that these
chartwise measures glue.  Explicitly, for $u$ a continuous function with
compact support and $\omega$ a smooth $(n,n)$-form, the assignment
\[
 u\longmapsto\int_Xu\,\omega
\]
is linear, and nonnegative when $u\ge0$ and $\omega$ is positive, by the last
clause of \Cref{prop:ch13-chart-integral}; the Riesz representation theorem
of \Cref{chap:topological-tools} then produces the measure, which is Radon
because $X$ is locally compact Hausdorff and the functional is locally
bounded by \Cref{prop:test-form-continuity-criterion}.  The measure does not
charge $\partial X$, by \Cref{cor:trop-boundary-negligible}.

Let $K'/K$ be a complete extension.  Refine a chart after base change so
that every target facet maps onto a facet of the original tropicalization.
Conservation of degree for graded reductions says that the sum of the
residue degrees and component lengths over a source facet equals its
original calibration.  Applying the tropical projection formula cellwise
therefore identifies the pushforward of the base-changed measure with the
original measure.  This proves scalar-extension compatibility.  The graded
degree statement and its non-strict-field formulation are
\cite[Propositions~15.3.1--15.3.2]{CLD25}.  In the language of
\Cref{chap:analytic-tropicalization} the mechanism is
\Cref{prop:trop-calibration-base-change}: the base-change morphism
$p:X_{K'}\to X$ satisfies $p_*\mu_{\Sigma_{K'}}=\mu_\Sigma$ on every
polyhedral subspace, and the tropical targets are literally unchanged, so
pairing with a form gives $p_*\bigl((p^*\omega)_{\mathrm{Rad}}\bigr)
=\omega_{\mathrm{Rad}}$.
\end{proof}

\begin{corollary}
\label{cor:ch13-measure}
Let $X$ be topologically separated of pure dimension $n$ and let
$\omega\in A^{n,n}_{X,G}(X)$ be a $G$-form with locally integrable
coefficients.  Then there is a unique Radon measure $\omega_{\mathrm{Rad}}$
on $X$ with
\[
 \int_Xu\,\omega_{\mathrm{Rad}}=\int_Xu\,\omega
\]
for every smooth compactly supported $u$; it is positive when $\omega$ is
positive, its total mass on a compact analytic domain $V$ tropicalized by $f$
is computed by \eqref{eq:ch13-chart-integral}, and it gives no mass to
$\partial X$ nor to any weakly analytic subset of dimension $<n$.
\end{corollary}

\begin{proof}
Uniqueness holds because smooth compactly supported functions are dense in
$C_c(X)$ for a locally holomorphically separated space, by the
Stone--Weierstrass theorem \cite[Theorem~15.2.6]{CLD25}; existence and
positivity are contained in the proof of
\Cref{thm:cld-analytic-integration}.  For the last clause,
\Cref{prop:trop-weakly-analytic-image} shows that such a subset has tropical
image of dimension $<n$ in every chart, and a subset of dimension $<n$ is
$\omega_D$-null in each $n$-cell; the boundary case is
\Cref{cor:trop-boundary-negligible}.
\end{proof}

The last compatibility of the integral is with finite morphisms.  It is the
analytic counterpart of the projection formula of
\Cref{cor:ch11-projection-formula} and it is where the calibrated degree of
\Cref{sec:trop-graded} appears in analysis rather than in combinatorics.

\begin{proposition}
\label{prop:ch13-finite-integral}
Let $p:Y\to X$ be a morphism of topologically separated $K$-analytic spaces
of level $n$, finite above the complement of a Zariski-closed subset of $X$
of dimension $<n$, and let $\mathrm{dc}(p)$ be the calibrated degree of
\Cref{thm:trop-finite-pushforward}.  Then for every $\omega\in
A^{n,n}_{X,G}(X)$ with compact support and locally integrable coefficients,
$p^*\omega$ is again such a form on $Y$ and
\[
 p_*\bigl((p^*\omega)_{\mathrm{Rad}}\bigr)
 =\mathrm{dc}(p)\,\omega_{\mathrm{Rad}} .
\]
In particular, if $p$ is finite and flat of degree $d$ then
$\int_Yp^*\omega=d\int_X\omega$.
\end{proposition}

\begin{proof}
Both sides are Radon measures, so it suffices to test against smooth
compactly supported functions and, by a partition of unity, to work on a
compact analytic domain $V\subset X$ carrying a moment map $f$ with
$\Sigma\subset\Sigma_f$, replacing $Y$ by a compact analytic neighbourhood
$W$ of $p^{-1}(\Sigma)$.  Set $\Sigma=\Sigma_f^{(n)}\cap V$ and
$\mathrm T=p^{-1}(\Sigma)$.  Then $\mathrm T$ is a polyhedral subspace of $W$
with canonical calibration $\mu_{\mathrm T}$, and
\Cref{thm:trop-finite-pushforward} gives
$p_*\mu_{\mathrm T}=\mathrm{dc}(p)\,\mu_\Sigma$.  Now
$p^*\omega|_{\mathrm T}=(p|_{\mathrm T})^*(\omega|_\Sigma)$, because
restriction of forms to polyhedral subspaces commutes with pullback along a
morphism carrying one into the other.  Pairing the two sides of the
calibration identity with $\omega|_\Sigma$ and integrating, using the
change-of-variables formula \Cref{prop:ch11-pushforward-integral} for the
piecewise immersion $p|_{\mathrm T}$, gives the asserted equality of
measures.  The flat case is the statement that
$\mathrm{dc}_x(p)=d$ for every $x$
(\Cref{prop:trop-calibrated-degree}\,\cref{itm:trop-dc:flat}).  The general
statement, including the case where $p$ is not flat, is
\cite[Proposition~15.4.3]{CLD25}.
\end{proof}

\begin{theorem}\label{thm:analytic-is-tropical-space}
Let $X$ be topologically separated and pure of dimension $n$.  With its
analytic boundary and canonical calibration, $X$ is a tropical space of
level $n$ in the sense of CLD.  Its analytic interior is harmonic.
\end{theorem}
\begin{proof}
Choose a point of $X$ and a moment map whose tropicalization has dimension
$n$ near that point.  After an integer-affine generic projection, we may
assume that the tropicalization is generically finite over an open subset of
$\mathbf R^n$.  At a codimension-one cell $\tau$, take a small transversal
interval.  The corresponding analytic restriction is a finite morphism of
one-dimensional affinoid germs.  Generic flatness makes its degree constant
on each connected side.  Conservation of degree at the point over $\tau$
therefore gives
\[
              \sum_{\sigma\supset\tau}
                 m_\sigma u_{\sigma/\tau}=0
                 \quad\text{in }N/N_\tau,
\]
which is the balancing equation.  Equivalently, pairing with every character
reduces the identity to the degree-zero divisor formula on the transversal
analytic curve, as in the proof of
\Cref{thm:analytic-tropical-balancing}.  Product moment maps show that the
result is independent of the chart and that the canonical calibrations agree
under transitions.  Thus the analytic interior has the required harmonic
balanced tropical charts.

For a point on the analytic boundary, the same local finite map is made over
a half-polyhedron.  A face contributes to the calibrated tropical boundary
exactly when the corresponding affinoid germ has nonempty relative analytic
boundary; interior faces occur in opposite pairs and cancel.  This identifies
the parlinear boundary with the analytic boundary and completes the tropical
space axioms.  The finite-flat degree argument and the boundary comparison
are \cite[Theorems~15.6.1 and 15.7.1]{CLD25}.
\end{proof}

\begin{remark}[What \Cref{thm:analytic-is-tropical-space} actually supplies]
\label{rem:ch13-tropical-structure-contents}
The statement bundles five separate properties, and it is worth listing them
because different later results use different ones.
\begin{enumerate}
\item\label{itm:ch13-ts:charts} \emph{Charts.} Every point has a
  neighbourhood carried by a tropical chart, so forms and calibrations are
  defined; this is \Cref{lem:trop-charts-basis} and it needs goodness, which
  is automatic in the interior.
\item\label{itm:ch13-ts:calib} \emph{A canonical calibration.} Provided by
  \Cref{cor:trop-calibration-glue}, and it is the unique piece of data that
  is genuinely global.
\item\label{itm:ch13-ts:harm} \emph{Interior harmony.} The defect
  $\mathbf N(\tau)$ of \Cref{def:ch11-defect} vanishes at every
  codimension-one cell not meeting $\partial X$.  This is the substantive
  input, and the two proofs of it --- the curve-slope route of
  \Cref{thm:analytic-tropical-balancing} and the finite-flat degree count of
  \cite[Theorem~15.6.1]{CLD25} --- are compared in
  \Cref{rem:trop-balancing-two-routes}.
\item\label{itm:ch13-ts:bdry} \emph{Boundary identification.} The parlinear
  boundary of the calibrated structure is the analytic boundary $\partial X$.
\item\label{itm:ch13-ts:thin} \emph{Thinness of the boundary.} $\partial X$
  is weakly analytic of dimension $\le n-1$
  (\Cref{prop:trop-boundary-thin}), so it is negligible for $n$-dimensional
  integration but carries the $(n-1)$-dimensional boundary measure.
\end{enumerate}
Purity of dimension is not decoration: if $X$ has a nonempty compact analytic
domain of dimension $<n$, then that domain has empty local skeleton, its
constant tropicalization is covered by no skeleton, and $X$ is not rich in
the sense required.  This is \cite[Remark~15.7.3]{CLD25}, and it is the
reason every statement in \Cref{sec:ch13-integration,sec:ch13-stokes} says
\emph{purely} of dimension $n$.
\end{remark}

\section{Stokes' formula on analytic domains}
\label{sec:ch13-stokes}

We now have all the ingredients for the boundary formula.  The tropical
theorem is \Cref{thm:cld-global-stokes}; the analytic content is entirely in
\Cref{thm:analytic-is-tropical-space}, and specifically in items
\cref{itm:ch13-ts:harm} and \cref{itm:ch13-ts:bdry} of
\Cref{rem:ch13-tropical-structure-contents}.  What remains is to say exactly
what the right-hand side of the formula is, to record the hypotheses, and to
check the result on an example where every number can be computed.

\subsection{The boundary integral}
\label{sec:ch13-boundary-integral}

\begin{definition}[Boundary measure of an $(n-1,n)$-form]
\label{def:ch13-boundary-integral}
\index{boundary measure}
Let $X$ be topologically separated of pure dimension $n$, locally
holomorphically separated and paracompact, and let
$\omega\in A^{n-1,n}_{X,G}(X)$ have compact support and $C^1$ coefficients.
Choose a finite family of compact analytic domains $V_i$ covering
$\Supp\omega$, moment maps $f_i$ tropicalizing $\omega|_{V_i}$ with
representatives $\alpha_i$, cell decompositions $\mathscr D_i$ of
$\Pi_{f_i}$ adapted to the canonical calibration, and a smooth partition of
unity $(\rho_i)$ with $\sum_i\rho_i=1$ near $\Supp\omega$.  Enlarging each
$f_i$ to a product chart, we may and do assume that $f_i$ tropicalizes
$\rho_i$ as well as $\omega|_{V_i}$, so that $\rho_i\alpha_i$ is again a
superform on $\Pi_{f_i}$.  Set
\begin{equation}
\label{eq:ch13-boundary-integral}
 \int_X^{\partial}\omega
 =(-1)^{n-1}\sum_i\ \sum_{\tau\in(\mathscr D_i)_{n-1}}
   \int_\tau\bigl(\iota''_{\mathbf N(\tau)}(\rho_i\alpha_i)\bigr)\big|_\tau ,
\end{equation}
where $\mathbf N(\tau)\in N/N_\tau$ is the boundary defect
(\Cref{def:ch11-defect}) of the calibrated complex
$(\mathscr D_i,m_i)$ and $\iota''$ is contraction in the double-primed block,
taken in the block convention of
\Cref{def:ch09-contraction,def:ch11-contractions} without a traversal sign;
the prefactor $(-1)^{n-1}$ is the one of
\eqref{eq:superform-boundary-traces}.  For $\theta\in A^{n,n-1}_{X,G}(X)$ with
compact support and $C^1$ coefficients, make the same choices with
representatives $\beta_i$ of $\theta|_{V_i}$ in place of the $\alpha_i$, and
set
\[
 \int_X^{\partial}\theta
 =-\sum_i\ \sum_{\tau\in(\mathscr D_i)_{n-1}}
   \int_\tau\bigl(\iota'_{\mathbf N(\tau)}(\rho_i\beta_i)\bigr)\big|_\tau .
\]
\end{definition}

\begin{proposition}
\label{prop:ch13-boundary-welldefined}
The quantity \eqref{eq:ch13-boundary-integral} is independent of all the
choices, and
\[
 \int_X^{\partial}\omega=\int_X\dprime\omega .
\]
If the partition of unity is chosen with
$\Supp\rho_i\subset\Int(V_i)$, then moreover only those $\tau$ whose
preimage in $\Sigma_{f_i}^{(n)}$ meets $\partial X$ contribute, and in
particular $\int^\partial_X\omega=0$ whenever
$\Supp\omega\cap\partial X=\varnothing$.  The analogous statements hold for
$\theta$ with $\ddprime$.
\end{proposition}

\begin{proof}
Fix the data and apply \Cref{thm:ch11-defect-stokes} to the calibrated
complex $(\mathscr D_i,m_i)$ and to the compactly supported form
$\rho_i\alpha_i$; the theorem is applicable because
$(\mathscr D_i,m_i)$ is a pure $n$-dimensional weighted rational polyhedral
complex and $\rho_i\alpha_i$ has compact support in $|\mathscr D_i|$.  It
gives
\[
 \int_{(\mathscr D_i,m_i)}\dprime(\rho_i\alpha_i)
 =(-1)^{n-1}\sum_{\tau\in(\mathscr D_i)_{n-1}}
   \int_\tau\bigl(\iota''_{\mathbf N(\tau)}(\rho_i\alpha_i)\bigr)\big|_\tau .
\]
Summing over $i$, the left-hand side is
$\sum_i\int_X\dprime(\rho_i\omega)$ by \Cref{prop:ch13-chart-integral} and
the fact that $\dprime$ commutes with the chart representation
(\Cref{prop:cld-intrinsic-calculus}).  Now $\sum_i\rho_i\omega=\omega$ on all
of $X$, because the $\rho_i$ sum to $1$ near $\Supp\omega$ and each
$\rho_i\omega$ vanishes off $\Supp\omega$; applying the $\mathbf R$-linear
operator $\dprime$ to this finite sum gives
$\sum_i\dprime(\rho_i\omega)=\dprime\omega$, and hence
$\sum_i\int_X\dprime(\rho_i\omega)=\int_X\dprime\omega$.  The right-hand
side is \eqref{eq:ch13-boundary-integral}.  Hence
the two agree, and since the left-hand side is independent of all choices by
\Cref{thm:cld-analytic-integration}, so is the right-hand side.

For the localization on $\partial X$: by
\Cref{rem:ch13-tropical-structure-contents}\,\cref{itm:ch13-ts:harm}, the
defect $\mathbf N(\tau)$ vanishes for every codimension-one cell $\tau$
whose preimage in $\Sigma^{(n)}_{f_i}$ lies in $\Int(V_i)$, so those terms
drop out of \eqref{eq:ch13-boundary-integral}; and the contraction
$\iota''_{0}=0$.  It is $\Int(V_i)$ and not $\Int(X)$ that appears here:
$\mathbf N(\tau)$ is the defect of the calibrated complex of the compact
analytic domain $V_i$, and $\partial V_i$ contains the relative boundary
$\partial(V_i/X)$, which lies inside $\Int(X)$.  Cells over that relative
boundary carry nonzero defect --- the two Shilov defects $\pm1$ of
\Cref{ex:ch13-stokes-annulus} are exactly of this kind --- and cancel in the
total rather than dropping out term by term.  The hypothesis
$\Supp\rho_i\subset\Int(V_i)$ removes them from the sum in the first place;
with it, if $\Supp\omega$ misses $\partial X$ altogether then every
surviving $\tau$ is disjoint from the support and the whole sum vanishes.
Without it, the identity $\int_X^\partial\omega=\int_X\dprime\omega$ still
holds, since it is proved above independently.
\end{proof}

\begin{theorem}
\label{thm:ch13-stokes}
\index{Stokes formula@Stokes' formula!on an analytic space}
Let $K$ be a complete non-Archimedean field and let $X$ be a topologically
separated $K$-analytic space, purely of dimension $n$, locally
holomorphically separated and paracompact, so that the partitions of unity of
\Cref{thm:cld-smooth-partition} are available.  Let
$\omega\in A^{n-1,n}_{X,G}(X)$ be a $G$-form with $C^1$ coefficients and
compact support, and let $\theta\in A^{n,n-1}_{X,G}(X)$ be one of complementary
type.  Then
\[
 \int_X\dprime\omega=\int_X^{\partial}\omega,
 \qquad
 \int_X\ddprime\theta=\int_X^{\partial}\theta ,
\]
with the boundary integrals of \Cref{def:ch13-boundary-integral}.  In
particular:
\begin{enumerate}
\item\label{itm:ch13-stokes:bdryless} if $\partial X=\varnothing$, then
  $\int_X\dprime\omega=0$ and $\int_X\ddprime\theta=0$;
\item\label{itm:ch13-stokes:interior} the same conclusion holds without
  assuming $\partial X=\varnothing$ provided
  $\Supp\omega\subset\Int(X)$, respectively
  $\Supp\theta\subset\Int(X)$;
\item\label{itm:ch13-stokes:domain} if $V$ is an analytic domain of a
  boundaryless $X$ and $\omega$ is supported in a compact subset of $V$, then
  \[
   \int_V\dprime\omega=\int^{\partial}_{V}\omega,
  \]
  the boundary measure of $V$ being carried by the relative boundary
  $\partial(V/X)$.
\end{enumerate}
\end{theorem}

\begin{proof}
The two displayed identities are \Cref{prop:ch13-boundary-welldefined}.
Assertion \cref{itm:ch13-stokes:bdryless} follows because a boundaryless $X$
has $\mathbf N(\tau)=0$ everywhere by interior harmony, so all boundary terms
vanish; equivalently one may quote \Cref{thm:cld-global-stokes} for the
calibrated tropical space of \Cref{thm:analytic-is-tropical-space}.
Assertion \cref{itm:ch13-stokes:interior} is the last clause of
\Cref{prop:ch13-boundary-welldefined}.

For \cref{itm:ch13-stokes:domain}, note first that $V$ carries two a priori
different structures of calibrated tropical space of level $n$: one as an
analytic space in its own right, one as a domain of the tropical space $X$.
These coincide, because the boundaries agree:
\[
 \partial V=(\partial X\cap V)\cup\partial(V/X)=\partial(V/X)
\]
when $\partial X=\varnothing$, and the canonical calibrations agree by
\Cref{cor:trop-calibration-restriction}.  Applying the first identity on
$V$ therefore gives $\int_V\dprime\omega=\int^\partial_V\omega$, and
\Cref{prop:ch13-boundary-welldefined} localizes the right-hand side on
$\partial V=\partial(V/X)$.
\end{proof}

\begin{warning}{Every hypothesis in \Cref{thm:ch13-stokes} is used}
\emph{Purity of dimension} is needed for the calibrated structure to exist at
all (\Cref{rem:ch13-tropical-structure-contents}).  \emph{Compact support} is
needed twice: to make the sums finite, and because
\Cref{thm:ch11-defect-stokes} is false for forms with noncompact support ---
already on $\mathbf R$, taken as a single unbounded cell of weight one, the
form $\eta=\varphi(x)\,\ddprime x$ with $\varphi$ smooth, $\varphi\equiv0$
near $-\infty$ and $\varphi\equiv1$ near $+\infty$, has
$\int_{\mathbf R}\dprime\eta=\int_{\mathbf R}\varphi'=1$, while the complex
has no codimension-one cell and hence no boundary term to absorb it; the
correct statement requires either compact support or a limiting
hypothesis.
\emph{The $G$-form hypothesis} matters when $\partial X\ne\varnothing$:
boundary points of a compact analytic domain have no ordinary neighbourhood
in the domain, and an ordinary form would not see them.  \emph{Topological
separation} is what makes supports behave, and \emph{local holomorphic
separation together with paracompactness} is what supplies the partition of
unity used to define both sides (\Cref{thm:cld-smooth-partition}); without
them the identity survives only chartwise, as \Cref{sec:ch13-properness}
records.  What is \emph{not} needed is goodness of $X$, algebraic closedness
of $K$, or nontriviality of the valuation.
\end{warning}

\begin{corollary}\label{cor:analytic-stokes-green}
If $X$ is pure $n$-dimensional and boundaryless, then
\[
 \int_X\dprime\eta=0,
 \qquad \int_X\ddprime\theta=0
\]
for $\eta\in A_{X,c}^{n-1,n}$ and
$\theta\in A_{X,c}^{n,n-1}$.  Consequently, for closed complementary
$\gamma$ and smooth $u,v$ with compact intersection of supports,
\[
 \int_Xu\,\dprimeddprime v\wedge\gamma
 =\int_Xv\,\dprimeddprime u\wedge\gamma.
\]
\end{corollary}
\begin{proof}
Apply the CLD Stokes and Green formulae
\cite[Proposition~5.11.1]{CLD25} to the tropical space supplied by
\Cref{thm:analytic-is-tropical-space}.  For clarity, the Green
identity uses integration by parts twice: first move $\dprime$ from
$\ddprime v\wedge\gamma$ to $u$, and then move $\ddprime$ from $v$ to
$\dprime u\wedge\gamma$.  The signs cancel because $\gamma$ is closed.
The same calculation with $u$ and $v$ interchanged gives the displayed
symmetry.  Applying Stokes once to the difference would leave cross terms
and is not, by itself, a proof.
\end{proof}

The boundary version is obtained by keeping the boundary terms rather than
discarding them; it is the identity that produces Dirichlet problems and
capacities in \Cref{chap:cld-bedford-taylor}.

\begin{theorem}
\label{thm:ch13-green-boundary}
Let $X$ be topologically separated of pure dimension $n$, and let $\alpha$
and $\beta$ be symmetric $G$-forms with $C^2$ coefficients, of bidegrees
$(p,p)$ and $(q,q)$ with $p+q=n-1$, whose supports have compact
intersection.  Then
\[
 \int_X\bigl(\alpha\wedge\dprimeddprime\beta
              -\dprimeddprime\alpha\wedge\beta\bigr)
 =\int_X^{\partial}\bigl(\alpha\wedge\ddprime\beta
                          -\ddprime\alpha\wedge\beta\bigr).
\]
If $\partial X=\varnothing$, or if the intersection of the supports lies in
$\Int(X)$, the right-hand side vanishes.
\end{theorem}

\begin{proof}
Put $\omega=\alpha\wedge\ddprime\beta-\ddprime\alpha\wedge\beta$, a $G$-form
of bidegree $(n-1,n)$ with $C^1$ coefficients, whose support is contained in
$\Supp\alpha\cap\Supp\beta$ and hence compact.  The Leibniz rule
(\Cref{prop:superform-differential-identities}) gives
\[
 \dprime\omega
 =\dprime\alpha\wedge\ddprime\beta
  +(-1)^{2p}\alpha\wedge\dprimeddprime\beta
  -\dprimeddprime\alpha\wedge\beta
  -(-1)^{2p+1}\ddprime\alpha\wedge\dprime\beta,
\]
and the two mixed terms cancel by the identity
$\dprime\alpha\wedge\ddprime\beta=-\ddprime\alpha\wedge\dprime\beta$ for
symmetric forms of complementary bidegree, which is
\Cref{lem:ch11-mixed-cancellation}.  Hence
$\dprime\omega=\alpha\wedge\dprimeddprime\beta-\dprimeddprime\alpha\wedge\beta$
and the assertion is \Cref{thm:ch13-stokes} applied to $\omega$, together
with \Cref{prop:ch13-boundary-welldefined} for the vanishing statements.
\end{proof}

\subsection{Stokes on an analytic annulus, with numbers}
\label{sec:ch13-stokes-annulus}

The following example is the one-dimensional unit test.  Every ingredient of
\Cref{thm:ch13-stokes} --- the calibration, the outward normal, the two
signs, the boundary measure --- occurs, and each can be evaluated by hand.

\begin{example}[Stokes on a closed annulus in $\mathbf P^{1,\An}$]
\label{ex:ch13-stokes-annulus}
Let $K$ be complete and nontrivially valued, let $q\in K^\times$ with
$0<|q|<1$, and put $\lambda=v(q)=-\log|q|>0$.  Inside the boundaryless space
$X=\mathbf P^{1,\An}_K$ consider the compact analytic domain
\[
 V=\{x\in X:\ |q|\le|T(x)|\le1\},
\]
a closed annulus.  We take $n=1$ throughout, so $c_1=1$ and
$\tau_1=\dprime t\wedge\ddprime t$.

\emph{Chart, skeleton, calibration.}  The function $T$ is invertible on $V$,
so $f=T:V\to\mathbf G_m$ is a moment map, and
\[
 t=f_{\mathrm{trop}}=-\log|T|=v(T)
\]
takes values in $[0,\lambda]$.  By \Cref{ex:trop-torus-explicit} the local
skeleton is $\Sigma_f=\Sigma_f^{(1)}=\{\eta_t:0\le t\le\lambda\}$, the
segment of Gauss points, mapped by $f_{\mathrm{trop}}$ isomorphically onto
$\Pi_f=[0,\lambda]\subset\mathbf R$.  The map $f$ is injective on
$\Sigma_f$ and finite flat of degree one over each Gauss point, so the
canonical calibration is the constant primitive $1$-vector: in the notation
of \eqref{eq:trop-multiplicity-formula}, $\mathscr D=\{[0,\lambda]\}$ with
weight $m=1$, and the lattice measure $\omega_{[0,\lambda]}$ is $\rd t$.
The total lattice length of $\Pi_f$ is $\lambda=v(q)$, matching
\Cref{prop:ch08-annulus-skeleton} and the edge-length convention of
\Cref{ex:ch08-An-edge}.

\emph{Boundary.}  Since $X$ is boundaryless,
$\partial V=\partial(V/X)$ consists of the two Shilov points $\eta_0$ and
$\eta_\lambda$, corresponding to $t=0$ and $t=\lambda$.  In the tropical
picture, $\mathscr D_0=\{\{0\},\{\lambda\}\}$ and the boundary defects of
\Cref{def:ch11-defect} are
\[
 \mathbf N(\{\lambda\})=m\cdot n_{\sigma,\{\lambda\}}=+1,
 \qquad
 \mathbf N(\{0\})=m\cdot n_{\sigma,\{0\}}=-1,
\]
because the outward primitive normal of $\sigma=[0,\lambda]$ at its right
endpoint is $+1$ and at its left endpoint is $-1$.

\emph{The form and both sides.}  Take
\[
 \omega=f_{\mathrm{trop}}^*\bigl(\varphi(t)\,\ddprime t\bigr)
 \in A^{0,1}_{X,G}(V),
 \qquad \varphi\in C^\infty(\mathbf R).
\]
Then $\dprime\omega=f_{\mathrm{trop}}^*(\varphi'(t)\,\dprime t\wedge\ddprime t)$,
so by \eqref{eq:ch13-chart-integral}
\[
 \int_V\dprime\omega
 =1\cdot\int_0^\lambda\varphi'(t)\,\rd t
 =\varphi(\lambda)-\varphi(0).
\]
On the other side, $n=1$ makes the prefactor $(-1)^{n-1}=+1$, and
$\iota''_{c}\bigl(\varphi\,\ddprime t\bigr)=c\,\varphi$, so
\eqref{eq:ch13-boundary-integral} reads
\[
 \int_V^\partial\omega
 =\bigl(+1\bigr)\varphi(\lambda)+\bigl(-1\bigr)\varphi(0)
 =\varphi(\lambda)-\varphi(0),
\]
the lattice measure of a point being counting measure of mass one.  The two
sides agree, as \Cref{thm:ch13-stokes}\,\cref{itm:ch13-stokes:domain}
predicts, and the computation is \Cref{lem:ch11-segment-stokes} transported
through one chart.

\emph{Numbers.}  Let $K=\mathbf C_p$ with $|p|=p^{-1}$, take $q=p^{3}$, so
$\lambda=v(q)=3\log p$, and take $\varphi(t)=t^2$.  Then
\[
 \int_V\dprime\omega=\lambda^2=9(\log p)^2,
 \qquad
 \int_V^{\partial}\omega=\varphi(\lambda)-\varphi(0)=9(\log p)^2 .
\]
If instead we normalize the valuation by $v(p)=1$, that is, replace
$-\log|\cdot|$ by $\log_p|\cdot|^{-1}$, then $\lambda=3$ and both sides
equal $9$.  Both sides scale in the same way, as they must: rescaling the
valuation by $c>0$ replaces $t$ by $ct$ and hence $\varphi$ by
$\varphi(\cdot/c)$, so the boundary values are unchanged, and the interior
integral is unchanged because $\rd t$ scales by $c$ while $\varphi'$ scales
by $c^{-1}$.  This is the one-dimensional case of the statement that a
lattice-normalized integral depends on the integral affine structure and not
on any choice of unit.

\emph{Two failures.}  First, drop compact support: take $\varphi\equiv1$ on
all of $\mathbf R$ and regard $\omega$ as a form on the open annulus
$\mathring V=\{|q|<|T|<1\}$, which is boundaryless as an analytic space.
Then $\dprime\omega=0$, so $\int_{\mathring V}\dprime\omega=0$, whereas the
naive ``boundary'' contribution $\varphi(\lambda)-\varphi(0)$ is also $0$;
but for $\varphi(t)=t$ one gets $\int_{\mathring V}\dprime\omega
=\int_0^\lambda\rd t=\lambda\ne0$ even though $\mathring V$ has no boundary.
There is no contradiction with
\Cref{thm:ch13-stokes}\,\cref{itm:ch13-stokes:bdryless}: the form
$t\,\ddprime t$ does not have compact support on $\mathring V$, since its
tropical support is all of $(0,\lambda)$ whose closure meets the two removed
Shilov points.  Second, forget that $V$ is compact and try to work with
ordinary forms: the Shilov points $\eta_0,\eta_\lambda$ have no ordinary
neighbourhood inside $V$, and the ordinary sheaf $A^{0,1}_V$ does not see
them.  This is exactly why \Cref{thm:ch13-stokes} is stated for $G$-forms.
\end{example}

\begin{example}[The same annulus after a degree-$m$ cover: the index appears
exactly once]
\label{ex:ch13-stokes-annulus-weight}
Keep $V$ and $\omega=f_{\mathrm{trop}}^*(\varphi(t)\,\ddprime t)$ as in
\Cref{ex:ch13-stokes-annulus}, assume $K$ algebraically closed, let $m\ge1$
be an integer, and let
\[
 V'=\{x:\ |q|^{1/m}\le|S(x)|\le1\},
 \qquad
 p:V'\longrightarrow V,\quad T=S^m,
\]
a finite flat morphism of degree $m$.  We compute
$\int_{V'}p^*\dprime\omega$ in two different charts on $V'$ and check that
the answers agree and equal $m\int_V\dprime\omega$, as
\Cref{prop:ch13-finite-integral} requires.

\emph{Chart $S$.}  Here $t'=v(S)$ ranges over $[0,\lambda/m]$, the moment map
$S$ is injective on the skeleton of $V'$ and finite flat of degree $1$ over
each Gauss point, so $\Pi_S=[0,\lambda/m]$ carries weight $1$ and lattice
measure $\rd t'$.  Since $t=v(T)=m\,v(S)=mt'$,
\[
 p^*\dprime\omega
 =\varphi'(mt')\,\dprime(mt')\wedge\ddprime(mt')
 =m^2\varphi'(mt')\,\dprime t'\wedge\ddprime t',
\]
so by \eqref{eq:ch13-chart-integral}
\[
 \int_{V'}p^*\dprime\omega
 =1\cdot\int_0^{\lambda/m}m^2\varphi'(mt')\,\rd t'
 =m\int_0^\lambda\varphi'(t)\,\rd t
 =m\bigl(\varphi(\lambda)-\varphi(0)\bigr).
\]

\emph{Chart $T\circ p=S^m$.}  Here the tropical coordinate is
$u=v(S^m)=mt'$, ranging over $[0,\lambda]$; the moment map $S^m:V'\to
\mathbf G_m$ is finite flat of degree $m$, so
\Cref{thm:trop-multiplicity-formula} gives $\Pi=[0,\lambda]$ with weight
$m_D=m$ and lattice measure $\rd u$.  Now
$p^*\dprime\omega=\varphi'(u)\,\dprime u\wedge\ddprime u$, whence
\[
 \int_{V'}p^*\dprime\omega
 =m\int_0^\lambda\varphi'(u)\,\rd u
 =m\bigl(\varphi(\lambda)-\varphi(0)\bigr).
\]

The two computations agree, and the factor $m$ appears once in each: in the
first chart through the square of the linear part of $u\mapsto mu$ against
the single lattice measure $\rd t'$ on a segment of lattice length
$\lambda/m$; in the second chart through the calibration weight.  This is the
lattice-index bookkeeping of the example following
\Cref{lem:trop-refinement-calculus}, carried out on an integral rather than
on a weight.  Finally,
$\int_{V'}p^*\dprime\omega=m\int_V\dprime\omega$ is
\Cref{prop:ch13-finite-integral} for the flat morphism $p$, and by
\Cref{thm:ch13-stokes} both sides equal $m$ times a difference of boundary
values --- the boundary of $V'$ being again two Shilov points, now at
$t'=0$ and $t'=\lambda/m$.
\end{example}

\section{Comparison with forms on a skeleton}
\label{sec:ch13-skeleton}

A form on $X$ is, by construction, a chartwise object; the calibrated
integration theory says that in top bidegree only its restriction to local
skeleta matters.  This section makes that precise, in the two directions
that get used: a faithfulness statement, which says forms are determined by
their skeletal restrictions, and a comparison statement, which identifies
$A^{p,q}_X$ with forms on a skeleton in favourable cases.  The Tate curve is
computed completely.

\subsection{Restriction to a polyhedral subspace}

Let $\Sigma\subset X$ be a polyhedral subspace in the sense of
\Cref{cor:trop-calibration-glue}, and let $j:\Sigma\to X$ denote the
inclusion.  A chart datum $(f,P)$ on an analytic domain $V$ restricts to the
polyhedral map $f_{\mathrm{trop}}|_{\Sigma\cap V}$, whose image lies in $P$;
pulling back a superform along it gives a form on the polyhedral space
$\Sigma\cap V$ in the sense of \Cref{def:ch10-forms-on-support}.  Since the
construction is compatible with morphisms of chart data, it descends to a
restriction morphism
\[
 \omega\longmapsto\omega|_\Sigma,
 \qquad
 A^{p,q}_{X,G}\longrightarrow j_*A^{p,q}_\Sigma ,
\]
compatible with $\wedge$, $J$, $\dprime$, $\ddprime$ and with further
restriction to a polyhedral subspace of $\Sigma$.  This is the analytic
instance of \cite[\S3.3.13]{CLD25}.

\begin{proposition}
\label{prop:ch13-skeletal-faithful}
Let $X$ be of level $n$ and let $\omega\in A^{p,q}_{X,G}(V)$ with
$\max(p,q)=n$, where $V$ is a compact analytic domain tropicalized by a
moment map $f$.  Then
\[
 \omega=0
 \iff
 \omega|_{\Sigma_f^{(n)}}=0
 \iff
 \alpha|_{\Pi_f}=0,
\]
where $\omega=f_{\mathrm{trop}}^*\alpha$.  Consequently a form of bidegree $(p,q)$ with $\max(p,q)=n$ on $X$ is zero as
soon as its restriction to every top-dimensional local skeleton vanishes.
\end{proposition}

\begin{proof}
By \Cref{prop:ch13-nullity} we have $\omega=0$ if and only if
$\alpha|_{\Trop_f(V)}=0$.  Decompose $\Trop_f(V)$ into cells; a cell $C$ of
dimension $<n$ has $\dim N_{C,\mathbf R}<n=\max(p,q)$, so the condition
\eqref{eq:ch10-form-kernel} at points of $C$ is vacuous, since an
alternating form of degree $n$ vanishes on any $n$-tuple of vectors from a
space of dimension $<n$.  Hence $\alpha|_{\Trop_f(V)}=0$ if and only if
$\alpha$ vanishes on the $n$-dimensional cells, that is, on
$\Trop_f(V)^{(n)}=\Pi_f$ (\Cref{def:trop-tropical-polyhedron}, using that
$V$ is compact).  Finally
$f_{\mathrm{trop}}:\Sigma_f^{(n)}\to\Pi_f$ is a surjective piecewise
immersion, so $\alpha|_{\Pi_f}=0$ if and only if
$\omega|_{\Sigma_f^{(n)}}=(f_{\mathrm{trop}}|_{\Sigma_f^{(n)}})^*
(\alpha|_{\Pi_f})=0$, by
\Cref{lem:ch13-surjective-pullback-injective} applied cellwise.

For the global statement, a form vanishing on all local skeleta vanishes on
each member of a $G$-cover by compact analytic domains, hence vanishes.
\end{proof}

\begin{warning}{Faithfulness fails below the top bidegree}
\Cref{prop:ch13-skeletal-faithful} is stated for $\max(p,q)=n$ and is false
in general otherwise: in bidegree $(0,0)$ the top-dimensional local skeleta
of a single chart need not detect a function.  The Berkovich affine line is
misleading here.  There $\Sigma_T$ is the real line of Gauss points
$\eta_t$, and a smooth function, being constant on every disc, \emph{is}
determined by its restriction to $\Sigma_T$; but that is a one-dimensional
accident, and it is not what the proposition asserts.  In general one must use all local skeleta, of
all dimensions, and the correct statement in low bidegree is that a form is
determined by its restrictions to the loci $\Sigma_f$ of all moment maps.
The reason integration only needs the top-dimensional part is precisely the
dimension count in the proof above.
\end{warning}

\subsection{A skeleton that carries all the forms: the Tate curve}

\begin{example}[Forms on a Tate curve]
\label{ex:ch13-tate-forms}
Let $K$ be complete, algebraically closed and nontrivially valued, let
$q\in K^\times$ with $0<|q|<1$, put $\lambda=v(q)$, and let
\[
 E=\mathbf G_{m,K}^{\An}\big/q^{\mathbf Z}
\]
be the Tate curve, with quotient map $\pi:\mathbf G_m^{\An}\to E$ and
skeleton
\[
 \Sk(E)=\pi\bigl(S(\mathbf G_m)\bigr)\simeq\mathbf R/\lambda\mathbf Z,
\]
a circle of lattice circumference $\lambda$, as computed in
\Cref{ex:trop-tate-curve} and, from the model side, in
\Cref{ex:ch08-tate-curve}.  Write $\rho:E\to\Sk(E)$ for the retraction.  We
claim:
\[
 A^{p,q}_E(E)\ \simeq\ A^{p,q}\bigl(\mathbf R/\lambda\mathbf Z\bigr)
 \qquad\text{for all }p,q,
\]
the right-hand side being superforms on the circle with its integral affine
structure; explicitly $A^{0,0}=C^\infty(\mathbf R/\lambda\mathbf Z)$,
$A^{1,0}=C^\infty\cdot\dprime s$, $A^{0,1}=C^\infty\cdot\ddprime s$,
$A^{1,1}=C^\infty\cdot\dprime s\wedge\ddprime s$, and $A^{p,q}=0$ when
$\max(p,q)\ge2$.

\emph{Injectivity.}  $E$ is purely of dimension $1$ and $\Sk(E)$ is the
top-dimensional local skeleton of the two charts $f_1,f_2$ of
\Cref{ex:trop-tate-curve}.  For $\max(p,q)=1$ this is
\Cref{prop:ch13-skeletal-faithful}.  For $p=q=0$ we argue by differentiating.
Let $u$ be a smooth function on $E$ with $u|_{\Sk(E)}=0$.  Restriction to a
polyhedral subspace commutes with $\dprime$, so
\[
 (\dprime u)|_{\Sk(E)}=\dprime\bigl(u|_{\Sk(E)}\bigr)=0,
\]
and $\dprime u$ has bidegree $(1,0)$ with $\max(1,0)=1=\dim E$, so
$\dprime u=0$ by \Cref{prop:ch13-skeletal-faithful}.  Now let $V\subset E$ be
a connected open set on which $u=\varphi\circ f_{\mathrm{trop}}$ for a moment
map $f$ and a smooth $\varphi$.  Then $f_{\mathrm{trop}}^*(\dprime\varphi)=0$,
so $\dprime\varphi$ vanishes in $A^{1,0}(\Trop_f(V))$ by
\Cref{prop:ch13-nullity} applied on compact analytic domains exhausting $V$;
that is, all tangential derivatives of $\varphi$ along the cells of
$\Trop_f(V)$ vanish.  Since $\Trop_f(V)$ is the continuous image of a
connected set, it is connected, and a polyhedral set is path-connected by
paths lying in its cells, so $\varphi$ is constant on $\Trop_f(V)$ and $u$ is
constant on $V$.  Hence $u$ is locally constant on $E$; as $E$ is connected,
$u$ is constant, and $u|_{\Sk(E)}=0$ forces $u=0$.

\emph{Surjectivity.}  Let $\psi$ be smooth on $\mathbf R/\lambda\mathbf Z$
and lift it to a $\lambda$-periodic $\widetilde\psi$ on $\mathbf R$.  On the
chart domain $U_1=\pi(\{0<v(T)<\lambda\})$ the moment map $f_1$ has
$f_{1,\mathrm{trop}}=v(T)\circ s_1$, and on $U_2$ likewise, with
$f_{2,\mathrm{trop}}=f_{1,\mathrm{trop}}+\lambda$ on one of the two
components of $U_1\cap U_2$ (\Cref{ex:trop-tate-curve}).  Since
$\widetilde\psi$ is $\lambda$-periodic, the two local forms
$f_{j,\mathrm{trop}}^*\bigl(\text{expression in }\widetilde\psi\bigr)$ agree
on the overlap, so they glue to a global form on $E$ restricting to the
given one on $\Sk(E)$.  This is a second, concrete instance of refinement
independence: the transition is a translation, and translations act
trivially on the local skeleton and on its calibration by
\Cref{cor:trop-sigma-translation} and
\Cref{thm:trop-degree-translation}.

\emph{Vanishing above bidegree one} is
\Cref{prop:ch13-locality}\,\cref{itm:ch13-loc:dim}, since
$\dim E=1$.
\end{example}

\begin{example}[Integration over the skeleton of a Tate curve, against the
conventions of \Cref{chap:model-calculations}]
\label{ex:ch13-tate-integration}
Keep the notation of \Cref{ex:ch13-tate-forms}, and let $s$ be the affine
coordinate on $\Sk(E)=\mathbf R/\lambda\mathbf Z$ induced by $v(T)$.  Every
$(1,1)$-form on $E$ is $\omega=\varrho(s)\,\dprime s\wedge\ddprime s$ with
$\varrho$ smooth and $\lambda$-periodic, the canonical calibration is the
constant primitive $1$-vector (all local degrees are $1$, by
\Cref{ex:trop-tate-curve}), and \eqref{eq:ch13-chart-integral} gives
\begin{equation}
\label{eq:ch13-tate-integral}
 \int_E\omega=\int_0^\lambda\varrho(s)\,\rd s .
\end{equation}
Three consistency checks follow.

\emph{(i) Exact forms integrate to zero.}  $E$ is proper, hence
boundaryless, so \Cref{thm:ch13-stokes}\,\cref{itm:ch13-stokes:bdryless}
gives $\int_E\dprimeddprime u=0$ for every smooth $u$.  Directly: with
$u=\psi(s)$, \eqref{eq:ch13-tate-integral} gives
$\int_0^\lambda\psi''(s)\,\rd s=\psi'(\lambda)-\psi'(0)=0$ by periodicity.
The two computations agree, and the second exhibits the boundarylessness of
$E$ as the periodicity of the skeleton.

\emph{(ii) Total mass.}  Let $d\ge1$ and put $\varrho\equiv d/\lambda$, so
that $\int_E\omega=d$.  This is the CLD-side normalization of a form of total
mass $d$; on the model side, \Cref{ex:ch08-tate-curve} computes, for a line
bundle $L$ of degree $d$ on $E$ and a model $\cL$ with residual degrees
$d_1,\ldots,d_r$ on the components of a semistable model, the measure
$\sum_id_i\delta_{v_i}$, of total mass $\sum_id_i=d$.  The two agree in
total mass, as they must, and they differ in kind: a smooth CLD form produces
an \emph{absolutely continuous} measure with respect to the lattice length
$\rd s$, whereas a model metric produces an \emph{atomic} measure at the
divisorial points.  Neither is a smooth approximation of the other in the
sense of forms; the correct comparison passes through currents
(\Cref{chap:cld-currents}) and metrics
(\Cref{chap:cld-metrized-bundles,chap:compare-ma}).

\emph{(iii) Normalization of the circumference.}  The lattice length of
$\Sk(E)$ is $\lambda=v(q)$, not $1$ and not the number of components of a
model.  If $\lambda=r$ is a positive integer, the standard semistable model
has a cycle of $r$ projective lines and dual complex a circle subdivided into
$r$ unit edges (\Cref{ex:ch08-tate-curve}); the lattice length of each edge
is $1$, by the convention of \Cref{ex:ch08-An-edge} that a node with local
equation $xy=\pi^m$ contributes an edge of lattice length $m$.  Summing gives
$r=\lambda$, so the tropical and model normalizations agree with no
correction factor.  That they agree is a consequence of all the calibration
weights being $1$ here; \Cref{rem:ch08-base-change-lengths} records the
bookkeeping when a base change scales the valuation, and
\Cref{ex:trop-tate-curve} records the same statement on the tropical side.
Taking $\varrho\equiv d/\lambda$ and $\lambda=r=3$, $d=3$ gives
$\int_E\omega=3$, the same number as the total mass
$\delta_{v_1}+\delta_{v_2}+\delta_{v_3}$ computed in
\Cref{ex:ch08-tate-curve}.
\end{example}

\section{Where properness and boundarylessness are needed}
\label{sec:ch13-properness}

Several statements of this chapter and of the next have hypotheses that are
easy to lose track of, because in the applications one usually works with a
proper $X$, where all of them hold.  This section lists them and says which
conclusion each protects.

\paragraph{Boundarylessness is needed for Stokes without boundary terms, and
for nothing else.}
\Cref{thm:ch13-stokes} holds with no hypothesis on $\partial X$; it is only
the vanishing of the right-hand side that requires
$\partial X=\varnothing$, or the weaker condition
$\Supp\omega\subset\Int(X)$.  Consequently every ``integration by parts''
statement --- \Cref{cor:analytic-stokes-green},
\Cref{thm:ch13-green-boundary} with vanishing right-hand side, and the
symmetry of the Dirichlet pairing --- is a statement about the interior.  A
proper $K$-analytic space, for instance the analytification of a proper
scheme, is boundaryless, which is why the global theory of
\Cref{chap:cld-metrized-bundles} is free of boundary terms.

\paragraph{Properness of a moment map is needed for the global calibrated
tropical pushforward.}
\Cref{def:trop-tropical-polyhedron} requires
$f_{\mathrm{trop}}|_{\Sigma_f^{(n)}}$ to be topologically proper; otherwise
the global pushforward calibration used in
\eqref{eq:ch13-chart-integral} is not defined.  This is automatic on a
compact analytic domain, which is why every chartwise computation in this
chapter is made on compact domains and then assembled by a partition of
unity.  It is \emph{not} automatic on an open subset: on
$\mathbf G_m^{\An}$ the chart $T$ is skeletally proper, but on the open unit
disc minus the origin the same chart is not proper onto its image once one
allows the radius to tend to $0$, and integrals of non-compactly-supported
forms genuinely diverge.

\paragraph{Compactness of supports is needed for every integral formula.}
Both \Cref{thm:cld-analytic-integration} and \Cref{thm:ch13-stokes} fail
without it, as the second failure recorded in
\Cref{ex:ch13-stokes-annulus} shows.  Where a noncompact statement is wanted
--- for instance for the Monge--Amp\`ere measure of a psh function on a
noncompact space --- one either passes to Radon measures, as in
\Cref{cor:ch13-measure}, or imposes a growth condition, and the two are not
interchangeable.

\paragraph{Purity of dimension is needed for the calibrated structure.}
See \Cref{rem:ch13-tropical-structure-contents}.  On a space with components
of different dimensions there is no level-$n$ calibrated structure, and one
must work component by component; the integral of an $(n,n)$-form then
ignores the components of dimension $<n$, which is consistent but is a
convention, not a theorem.

\paragraph{Properness of a morphism is needed for pushforward, never for
pullback.}
\Cref{prop:cld-form-pullback} has no properness hypothesis: pullback of forms
along any morphism exists.  Pushforward does not exist for forms at all, and
for measures and currents it requires properness of the morphism on the
support: \Cref{prop:ch13-finite-integral} assumes finiteness away from a
thin set, and the fibre-integration and projection formulae of
\Cref{chap:cld-currents} assume relative properness of the support together
with density of the flat locus.  Dropping the flatness hypothesis makes the
projection formula false, the defect being measured by generic
multiplicities; the precise statement and counterexample are
\cite[Theorem~16.5.4 and Remark~16.5.6]{CLD25}.

\paragraph{Goodness, holomorphic separation and paracompactness are needed
for partitions of unity.}
This is the subject of the next section.  They are used to produce bump
functions and hence to globalize chartwise constructions; without them the
statements of this chapter remain true chartwise and on explicitly given
finite covers, and only the assembly step is lost.

\section{Partitions of unity and precise hypotheses}
\label{sec:ch13-partitions}

The topology of Berkovich spaces is not automatically second countable, and
smooth bump functions are not available without separation assumptions.
The version used throughout the book is the following.

\begin{theorem}\label{thm:cld-smooth-partition}
If $X$ is locally holomorphically separated and paracompact, every open cover
admits a subordinate smooth partition of unity.  Consequently the sheaves
$A_X^{p,q}$ are fine.
\end{theorem}
\begin{proof}
Local holomorphic separation gives, near every point, finitely many
invertible analytic functions whose logarithmic absolute values separate a
chosen compact neighborhood from a prescribed disjoint closed set.  Smooth
functions of these logarithms form an algebra containing constants and
separating the two compact sets.  Smooth Stone--Weierstrass approximation,
followed by a one-variable cutoff, produces a smooth function $\chi$ equal
to one on the smaller compact set and supported in the chosen open set.
Thus smooth compactly supported bump functions exist.

Paracompactness refines an arbitrary open cover to a locally finite family
of pairs $W_i\Subset V_i$, with the $W_i$ covering and each $V_i$ contained
in a member of the original cover.  Choose bumps $0\le\chi_i\le1$ satisfying
$\chi_i=1$ on $W_i$ and $\operatorname{supp}\chi_i\subset V_i$.  The sum
$s=\sum_i\chi_i$ is locally finite, smooth, and at least one.  Hence
\[
                         \rho_i=\chi_i/s
\]
is a locally finite subordinate smooth partition of unity.  Multiplication
by $(\rho_i)$ gives the usual sheaf-theoretic endomorphisms with locally
finite supports, so each $A_X^{p,q}$ is fine.  The separation and
approximation steps are
\cite[Proposition~15.2.8 and Corollary~15.2.10]{CLD25}.

Without the stated separation and paracompactness hypotheses, the argument
does not apply; later uses are then chartwise or use an explicitly given
finite cover.
\end{proof}

\begin{warning}{Fineness is not softness of the chart data}
\Cref{thm:cld-smooth-partition} makes $A^{p,q}_X$ a fine sheaf, so its higher
cohomology vanishes on a paracompact $X$.  This says nothing about the
cohomology of the complex $(A^{\bullet,q}_X,\dprime)$, which is the object of
interest and whose local exactness --- a Poincar\'e lemma --- is a separate
theorem, true for $\dprime$ and $\ddprime$ on smooth $K$-analytic spaces by
\cite{Jell16} and delicate in general.  Fineness of the terms of a complex
is a statement about partitions of unity; exactness is a statement about
the differentials.  We use only the former in this chapter, and return to the
latter in \Cref{chap:delta-forms}, where the enlargement of the sheaf by
harmonic tropicalizations is motivated precisely by the failure of the
smooth complex to compute the expected cohomology.
\end{warning}

\section{Dictionary established in this chapter}
\label{sec:ch13-dictionary}

The following correspondences are used without further comment from
\Cref{chap:cld-currents} onwards.
\[
\begin{array}{ccl}
\text{tropical chart datum }(f,P)
 &\longleftrightarrow&
 \text{object of }\mathscr T(V),\\[1mm]
\text{product moment map}
 &\longleftrightarrow&
 \text{common domination in }\mathscr T(V),\\[1mm]
f_{\mathrm{trop}}^*\alpha
 &\longleftrightarrow&
 \text{class in a filtered colimit, not a pullback},\\[1mm]
\alpha|_{\Trop_f(V)}=0
 &\longleftrightarrow&
 f_{\mathrm{trop}}^*\alpha=0\ (V\text{ compact}),\\[1mm]
\text{sheafification}
 &\longleftrightarrow&
 \text{passage from }A^{p,q}_{\mathrm{pf}}\text{ to }A^{p,q}_X,\\[1mm]
\text{cellwise kernel }K^{p,q}
 &\longleftrightarrow&
 \text{differential ideal},\\[1mm]
\text{morphism }g:Y\to X
 &\longleftrightarrow&
 g^*:A^{p,q}_X\to g_*A^{p,q}_Y,\\[1mm]
\text{local skeleton }\Sigma_f^{(n)}
 &\longleftrightarrow&
 \text{domain of integration},\\[1mm]
\text{canonical calibration}
 &\longleftrightarrow&
 \text{weights }m_D\text{ in }\eqref{eq:ch13-chart-integral},\\[1mm]
\text{lattice measure }\omega_D
 &\longleftrightarrow&
 \text{normalization of }\int_D,\\[1mm]
\text{boundary defect }\mathbf N(\tau)
 &\longleftrightarrow&
 \text{boundary measure }\textstyle\int^\partial,\\[1mm]
\text{analytic boundary }\partial X
 &\longleftrightarrow&
 \text{support of }\textstyle\int^\partial,\\[1mm]
\text{interior harmony}
 &\longleftrightarrow&
 \mathbf N(\tau)=0\text{ off }\partial X,\\[1mm]
\text{pure }n\text{-dimensional }X
 &\longrightarrow&
 \text{tropical space of level }n,\\[1mm]
\text{calibrated degree }\mathrm{dc}(p)
 &\longleftrightarrow&
 \text{factor in }p_*(p^*\omega)_{\mathrm{Rad}} .
\end{array}
\]
The second line is
\Cref{lem:ch13-charts-filtered}\,\cref{itm:ch13-cf:prod}; the fourth is
\Cref{prop:ch13-nullity}; the fifth and sixth are \Cref{thm:ch13-sheaf} and
\Cref{lem:ch13-kernel-differential}; the seventh is
\Cref{prop:cld-form-pullback}; the ninth and tenth are
\eqref{eq:ch13-chart-integral} and \Cref{def:ch11-lattice-measure}; the
eleventh through thirteenth are
\Cref{def:ch13-boundary-integral,prop:ch13-boundary-welldefined}; the
fourteenth is \Cref{thm:analytic-is-tropical-space}; and the last is
\Cref{prop:ch13-finite-integral}.

Five separations belong to the dictionary as much as the correspondences do,
and each is a place where a plausible-sounding statement is false.

First, \emph{$f_{\mathrm{trop}}^*$ is not a pullback and $g^*$ is}.  The
first is a colimit label with a computable kernel
(\Cref{prop:ch13-nullity}); the second is a genuine morphism of sheaves with
no kernel computation available in general
(\Cref{sec:ch13-what-is-pullback}).

Second, \emph{a form is not determined by pointwise values}.  It is
determined by its restrictions to local skeleta, and in top bidegree by the
top-dimensional ones only (\Cref{prop:ch13-skeletal-faithful}).  Below the
top bidegree, restriction to the top-dimensional skeleta loses information.

Third, \emph{ordinary forms are not $G$-forms}.  The comparison map is
injective and not surjective
(\Cref{thm:ch13-sheaf}\,\cref{itm:ch13-sheaf:GvsO}), and the difference is
carried by the boundary of a compact analytic domain.  All boundary formulae
are stated for $G$-forms.

Fourth, \emph{fineness is not exactness}.  The sheaves $A^{p,q}_X$ are fine
on a paracompact locally holomorphically separated space, but the complexes
they form need not be locally exact; see the warning above.

Fifth, \emph{the measure of a smooth form is not the measure of a model
metric}.  On the Tate curve the first is absolutely continuous on the
skeleton and the second is atomic
(\Cref{ex:ch13-tate-integration}); they agree in total mass and nowhere else.
The comparison of the two is the subject of
\Cref{chap:cld-metrized-bundles,chap:compare-ma} and requires currents.

\section*{Source notes}

The intrinsic analytic theory in this chapter follows
\cite[Chapters~15--16]{CLD25}; the affine superform calculus is due to
Lagerberg and was reconstructed in \Cref{chap:superforms}.  The
distinction between ordinary and $G$-forms, and between the ordinary and
$G$-topologies, is retained because it becomes essential for currents,
boundary terms, and fibre integration.

\paragraph{Provenance of the individual statements.}
The description of a form by moment maps and superforms is
\cite[Proposition~15.1.1 and Corollary~15.1.2]{CLD25}, which is our
\Cref{prop:cld-local-normal-form}.  Invariance of ordinary forms under
reduction is \cite[Proposition~15.1.3(b) and the form variant immediately
following it]{CLD25}; part~(a) and its corresponding variant govern
$G$-forms.  The filtered category of tropical
charts and the associated presheaf are \cite[\S3.3]{CLD25}; the nullity
criterion \Cref{prop:ch13-nullity} is their Propositions~3.4.1 and~3.4.4, and
the failure of separatedness recorded in the warning after it is their
Remark~3.4.3.  Our proof differs in one respect worth naming: Chambert-Loir
and Ducros deduce the criterion from the existence of piecewise affine
sections of a surjective affine map of polyhedra, whereas
\Cref{lem:ch13-surjective-pullback-injective} avoids sections entirely and
argues pointwise with the tangent spaces of a compatible decomposition.  The
two routes give the same theorem; the present one is chosen because the
same tangent-space computation reappears in
\Cref{lem:ch13-kernel-differential} and in
\Cref{prop:ch13-skeletal-faithful}, so the reader meets it once.

The vanishing of a form off the locus $d_K\ge\max(p,q)$ is
\cite[Proposition~15.2.1]{CLD25}; our \Cref{prop:ch13-locality} derives it
from the dimension chain of \Cref{chap:analytic-tropicalization} instead.
Holomorphic separation, the Stone--Weierstrass theorem for smooth functions,
the existence of smooth bump functions and of partitions of unity, and the
fineness of $A^{p,q}_X$ are \cite[\S15.2]{CLD25}, quoted in
\Cref{thm:cld-smooth-partition}.  Compatibility of integration with
extension of scalars is \cite[Propositions~15.3.1--15.3.2]{CLD25}, resting on
the base-change compatibility of the canonical calibration proved in
\Cref{prop:trop-calibration-base-change}.  Integration against a finite
morphism and the calibrated degree are \cite[\S15.4]{CLD25}, our
\Cref{prop:ch13-finite-integral}.  Interior harmony is
\cite[Theorem~15.6.1]{CLD25} and the identification of an analytic space with
a tropical space of level $n$ is \cite[Theorem~15.7.1]{CLD25}; both are quoted
in \Cref{thm:analytic-is-tropical-space}, whose statement and proof are those
of the first draft of this chapter and are not modified here.  The remark
that purity of dimension is necessary is their Remark~15.7.3.  The Stokes and
Green formulae for a calibrated $G$-tropical space are
\cite[Proposition~5.11.1]{CLD25}; our \Cref{thm:ch13-stokes} is that
statement transported through \Cref{thm:analytic-is-tropical-space}, with the
boundary measure defined by the defect formula of
\Cref{thm:ch11-defect-stokes} rather than by the abstract
$(n-1)$-dimensional volume forms of \cite[\S\S4.6, 5.9]{CLD25}.  The two
definitions agree, and the elementary one has the advantage that
\Cref{ex:ch13-stokes-annulus} can be evaluated without any further theory.
The projection formula and its flatness hypothesis, which we quote only to
delimit \Cref{sec:ch13-properness}, are \cite[Theorem~16.5.4 and
Remark~16.5.6]{CLD25}.

\paragraph{Two lineages, again.}
As in \Cref{chap:analytic-tropicalization}, the material has two ancestries.
The polyhedral one is Lagerberg's superform calculus, recalled in
\cite[\S1]{CLD25} and developed here in \Cref{chap:superforms}; the
convolution and regularization formalism of that calculus is what makes the
smooth theory usable, and its supercurrent counterpart is developed
independently by Berndtsson \cite{Ber19Super}, whose metric-dependent
supercurrent of a real submanifold is a different object from a calibrated
tropical current and is kept apart in \Cref{chap:cld-currents}.  The
analytic one is Berkovich's theory \cite{Berk12} together with Temkin's
graded reduction \cite{Temkin2004,Tem16}, which is what makes the local
skeleta of \Cref{chap:analytic-tropicalization} --- and hence the domains of
integration of this chapter --- intrinsic.  The tropical conventions,
including the lattice-index normalization used in
\Cref{ex:ch13-stokes-annulus-weight}, follow Mikhalkin--Rau
\cite[\S\S4.2, 4.6]{MR18Tropical}.

\paragraph{Later developments and their relation to this chapter.}
The sheaves constructed here are the smallest reasonable category of real
forms on a Berkovich space, and three enlargements are known.
Gubler--K\"unnemann delta-forms \cite{GK17,GK19} adjoin the curvature of
piecewise affine metrics and make tropical intersection an honest product;
they are the subject of \Cref{chap:delta-forms}.  Gubler--Jell--Rabinoff
forms based on harmonic tropicalizations \cite{GJR26} enlarge the supply of
charts by allowing harmonic rather than merely invertible coordinate
functions, with better cohomological behaviour; the Poincar\'e lemma for the
smooth sheaves, in the cases where it holds, is \cite{Jell16}.  The
curve-test theory of Gubler--Rabinoff \cite{GR25PL,GR25Classical} provides a
local criterion for semipositivity of piecewise linear functions which is
compared with the CLD notion only under stated hypotheses.  None of these is
used in the present chapter, and the comparison theorems are deferred to
\Cref{chap:delta-forms} so that the boundaries between the theories remain
visible.  For the identification of the CLD operator on a curve skeleton with
the classical Laplacian of \Cref{chap:model-calculations}, see
\cite{Wanner19}; the sign discrepancy there is the outward-versus-outgoing
convention discussed in \Cref{chap:tropical-integration} and not a
disagreement between sources.

\paragraph{What is not proved here.}
Three natural questions are outside the scope of this chapter and are
recorded so that they are not mistaken for consequences of what is proved.
The local exactness of $(A^{\bullet,\bullet}_X,\dprime,\ddprime)$ is not
established; the relation between $H^{p,q}$ of these sheaves and any
cohomology of $X$ is not discussed; and the surjectivity of the comparison
map $A^{p,q}_X\to A^{p,q}_{X,G}$, which fails, is not repaired by passing to
a larger class of charts.  Each of these is addressed, in the cases where an
answer is known, in \Cref{chap:delta-forms}.
